\section{Old Chinese vowels}

\paragraph{}
\label{par:old-chinese-vowels}
In the practice of all investigators the reconstruction of Old
Chinese rimes takes as a point of departure the internal reconstruction
of Middle Chinese rimes. To facilitate a linear discussion here, it is
convenient to focus on type A rimes with velar codas (§§-). The
hypotheses generated by considering the type A rimes with velar codas
can then be tested and refined with reference to the type A rimes with
dental codas (§§-). Together the consideration of distributional
relationships in type A syllables with velar and dental codas generates
a five vowel theory of Old Chinese which can next be tested against the
distribution of type B syllables; the type B rimes require the
postulation of a sixth vowel (§).

\begin{figure}
    \centering
\begin{tabular}{|l|l|l|l|l|l|}
\hline
& & P- &
T- &
Ts- &
K- \\ \hline
{\emph{-ang}} &
唐一開 &
{\emph{Pang}} &
{\emph{Tang}} &
{\emph{Tsang}} &
{\emph{Kang}} \\ \hline
{\emph{-wang}} &
唐一合 &
--- &
--- &
--- &
{\emph{Kwang}} \\ \hline
{\emph{-aeng}} &
庚二開 &
{\emph{Paeng}} &
{\emph{Taeng}} &
--- &
{\emph{Kaeng}} \\ \hline
{\emph{-waeng}} &
庚二合 &
--- &
--- &
--- &
{\emph{Kwaeng}} \\ \hline
{\emph{-uwng}} &
東一 &
{\emph{Puwng}} &
{\emph{Tuwng}} &
{\emph{Tsuwng}} &
{\emph{Kuwng}} \\ \hline
{\emph{-eng}} &
青四開 &
{\emph{Peng}} &
{\emph{Teng}} &
{\emph{Tseng}} &
{\emph{Keng}} \\ \hline
{\emph{-weng}} &
青四合 &
--- &
--- &
--- &
{\emph{Kweng}} \\ \hline
{\emph{-eang}} &
耕二開 &
{\emph{Peang}} &
--- &
{\emph{Tseang}} &
{\emph{Keang}} \\ \hline
{\emph{-weang}} &
耕二合 &
--- &
--- &
--- &
--- \\ \hline
{\emph{-ong}} &
登一開 &
{\emph{Pong}} &
{\emph{Tong}} &
{\emph{Tsong}} &
{\emph{Kong}} \\ \hline
{\emph{-wong}} &
登一合 &
--- &
--- &
--- &
{\emph{Kwong}} \\ \hline
{\emph{-owng}} &
冬一 & 
{\emph{Powng}} &
{\emph{Towng}} &
{\emph{Tsowng}} &
{\emph{Kowng}}
\\ \hline
\end{tabular}

    \caption{Middle Chinese division-i/iv syllable types in final -\textit{ng} 
    \citep[198]{Baxter2014OldChinese}}
    \label{fig:MChi-i/iv-final-ng}
\end{figure}


\begin{figure}
    \centering

\begin{tabular}{|l|l|l|l|l|l|l|}
 \hline
&&
P- &
T- &
Ts- &
K- &
K{\sil{ʷ}}- \\ \hline
{\emph{-ang}} &
唐一 &
{\emph{Pang}} &
{\emph{Tang}} &
{\emph{Tsang}} &
{\emph{Kang}} &
{\emph{K{\sil{ʷ}}ang}} \\ \hline
{\emph{-aeng}} &
庚二 &
{\emph{Paeng}} &
{\emph{Taeng}} &
--- &
{\emph{Kaeng}} &
{\emph{K{\sil{ʷ}}aeng}} \\ \hline
{\emph{-uwng}} &
東一 &
{\emph{Puwng}} &
{\emph{Tuwng}} &
{\emph{Tsuwng}} &
{\emph{Kuwng}} &
--- \\ \hline
{\emph{-eng}} &
青四 &
{\emph{Peng}} &
{\emph{Teng}} &
{\emph{Tseng}} &
{\emph{Keng}} &
{\emph{K{\sil{ʷ}}eng}} \\ \hline
{\emph{-eang}} &
耕二 &
{\emph{Peang}} &
--- &
{\emph{Tseang}} &
{\emph{Keang}} &
--- \\ \hline
{\emph{-ong}} &
登一 &
{\emph{Pong}} &
{\emph{Tong}} &
{\emph{Tsong}} &
{\emph{Kong}} &
{\emph{K{\sil{ʷ}}ong}} \\ \hline
{\emph{-owng}} &
冬一  &
{\emph{Powng}} &
{\emph{Towng}} &
{\emph{Tsowng}} &
{\emph{Kowng}} &
--- \\ \hline
\end{tabular}
    \caption{Middle Chinese division-i/iv syllable types in final -ng with the supposition of labiovelar initials}
    \label{fig:MChi-div-ng-with-labiovelars}
\end{figure}


\paragraph{Asymmetries in type A rimes with velar codas.}
presents the type A Middle Chinese rimes with final -\emph{ng}; their
distribution includes a number of unattractive asymmetries. In
particular the distribution of -\emph{w}- (cf. §) is quite irregular as
are the distributions of the vowels -\emph{ae}- and -\emph{ea}- (cf. §).

\paragraph{Origin of some -w- (合口
 \emph{\textbf{hékŏu}}\textbf{) rimes in labiovelar initials.}} The
rimes -\emph{wang}, -\emph{waeng}, -\emph{weng}, and -\emph{wong} occur
only with velar initials (牙 \emph{K}-). As proposed in § , this
dis­tri­bu­tion suggests the analysis of the -\emph{w}- element as part
of the initial rather than in the rime., and thereby to propose a series
of labiovelar initials. Analyzing the rhymes -\emph{wang},
-\emph{waeng}, -\emph{weng}, and -\emph{wong} as -\emph{ang},
-\emph{aeng}, -\emph{eng}, and -\emph{ong} after labiovelars allows to
be reformulated as . Some of the remaining gaps in this table are
phonetically well mo­ti­vated. In particular, the lack of syllables like
*k{\sil{ʷ}}uwng and *k{\sil{ʷ}}owng would have involved too much labialness in one
syllables. The lack of *k{\sil{ʷ}}eang is best treated as an accidental gap,
since \emph{hweang} (嶸) and \emph{kweak} (膕, 聝, etc.) do occur.


\paragraph{The `r-hypothesis': origin of the
vowels -ae- and -ea-.} Bearing in mind that the reconstruction of
Middle Chinese retroflex (知 \emph{Tr}-) and retroflex affricate
(\emph{zhuāngzǔ~}莊組 \emph{~}Tsr-) initials as deriving respectively
from Old Chinese dentals and affricates before medial *-r- is based on
the complementary distribution of type A division-ii vowels -\emph{ae}-
and -\emph{ea}- (division-ii) with dental initials (\emph{shétóu} 舌頭
 \emph{T}-) (\cite[203-206]{Baxter1992}, cf. §), the conjecture arises that
medial *-r- is also responsible for the emergence of these two vowels.

§a. Consideration of type B syllables supports the conjecture that
-\emph{aeng} originates from *‑aŋ. There are type B rimes with
-\emph{aeng}- after grave initials (\emph{kj}-,  \emph{pj}-, etc.), e.g.
京  \emph{kjaeng}, 兵  \emph{pjaeng}, but there are no type B rimes with
-\emph{aeng} after acute initials (\emph{trj}-,  \emph{tsy}- \textless{}
*t- etc.), i.e. syllables such as *trjaeng and *tsyaeng do not exist
\citep[272]{Baxter1992}. If one suggests that type B syllables in
-\emph{aeng} derive from *-raŋ, moving the distinction between
 \emph{kjang} and \emph{kjaeng} from the vowel to the medial makes
 \emph{kjaeng} \textless{} *kraŋ and \emph{kjang} \textless{} *kaŋ
parallel to \emph{trjang} \textless{} *traŋ and \emph{tsyang}
\textless{} *taŋ. If -\emph{aeng} derives from *-raŋ in type B
syllables, it is reasonable to suggest that -\emph{aeng} in type A
syllables also derives from *-raŋ. This internal reconstruction predicts
that words in Middle Chinese -\emph{ang} and -\emph{aeng} should rhyme
with each other in Old Chinese poetry. The rhymes of the 詩經
 \emph{Shījīng} bear out this prediction. Already the first few poems
offer the rhyme sequences 莫濩綌 \emph{mak hwak khjaek} (Ode 2.2) and
岡黃觥傷 \emph{kang hwang kwaeng syang} (Ode 3.3).\footnote{An
  exploration of rhyme behavior also reveals that Old Chinese *-raŋ is
  not the only origin of Middle Chinese -\emph{aeng}. In particular, in
  type B syllables -\emph{aeng} also descends from *-reŋ, e.g. 瑩
   \emph{hjwaeng} \textless{} *{\sil{[ɢ]ʷ}}reŋ `luster' (09-09k) rhymes with
  星  \emph{seng} \textless{} *s-ts{\sil{ʰ}}{\sil{ˤ}}eŋ `star' (09-25x) in Ode 55.2.}

\begin{figure}
    \centering

\begin{tabular}{|l|l|l|l|l|l|l|l|}
 \hline
OChi. &
MChi. &
&
P- &
T- &
Ts- &
K- &
Kʷ- \\ \hline
*-aŋ &
{\emph{-ang}} &
唐一 &
{\emph{Pang}} &
{\emph{Tang}} &
{\emph{Tsang}} &
{\emph{Kang}} &
{\emph{K{\sil{ʷ}}ang}} \\ \hline
*-oŋ &
{\emph{-uwng}} &
東一 &
{\emph{Puwng}} &
{\emph{Tuwng}} &
{\emph{Tsuwng}} &
{\emph{Kuwng}} &
--- \\ \hline
*-eŋ &
{\emph{-eng}} &
青四 &
{\emph{Peng}} &
{\emph{Teng}} &
{\emph{Tseng}} &
{\emph{Keng}} &
{\emph{K{\sil{ʷ}}eng}} \\ \hline
*-əŋ &
{\emph{-ong}} &
登一 &
{\emph{Pong}} &
{\emph{Tong}} &
{\emph{Tsong}} &
{\emph{Kong}} &
{\emph{K{\sil{ʷ}}ong}} \\ \hline
*-uŋ &
{\emph{-owng}} &
冬一  &
{\emph{Powng}} &
{\emph{Towng}} &
{\emph{Tsowng}} &
{\emph{Kowng}} &
--- \\ \hline
\end{tabular}
    \caption{Revised type A syllables with final *-ŋ}
    \label{fig:Revise-A-syll-with-ng}
\end{figure}

§b. Turning from the origin of -\emph{aeng} to the origin of
-\emph{eang}, the restriction of the latter to type A syllables already
points to its secondary origin. Since there is no sign of (former)
rounding in -\emph{eang}, it seem more reasonable to credit its origin
to *-eŋ or *-əŋ rather than to *oŋ or *-uŋ. In the 詩經  \emph{Shījīng}
syllables ending in -\emph{eang} rhyme both with -\emph{eng}, e.g. 丁城
 \emph{treang dzyeng} (Ode 7.1), and with -\emph{ong} {[}ʌŋ], e.g. 麥北
 \emph{meak pok} (Ode 48.2).\footnote{In fact I find only rhymes of
  -\emph{eak} and -\emph{ok} but none of -\emph{eang} and -\emph{ong}. }
It thus appears that both *-reŋ and *-rəŋ are origins of -\emph{eang}
and that these two originally distinct rimes merged on their way to
Middle Chinese.

If we rid of the rimes -\emph{aeng} and -\emph{eang}, the result is a
system of five vowels before velar codas. As a working hypothesis, we
may presume that the original nuclear vowels were *a, *e, *ə, *u, and *o
as shown in . In this analysis -\emph{ang} \textless{} *-aŋ and
-\emph{eng} \textless{} *-eŋ are retentions from the original system.
Recalling that Middle Chinese orthographic \textless{}o\textgreater{}
symbolizes a vowel similar to {[}ʌ], -\emph{ong} \textless{} *-əŋ is
also tantamount to a retention.

\paragraph{Asymmetries in type A rimes with dental codas.} The
search for asymmetries in type A syllables with velar codas provides a
five vowel theory of Old Chinese. To explore the ramifications of this
hypothesis further, discussion now turns to the type A rimes with dental
codas.

§a. The discussion of -\emph{ae}- and -\emph{ea}- in the case of the
velar final (cf. §) rimes applies \emph{mutatis mutandis} to the
appearance of these two vowels with dental finals. With the hypotheses
-\emph{ae}- \textless{} *-ra and -\emph{ea}- \textless{} *-re-, *‑rə-
one may for the present regard -\emph{aen}, -\emph{waen}, -\emph{ean},
and -\emph{wean} as deriving from -\emph{an}, -\emph{wan}, -\emph{en},
-\emph{on}, -\emph{wen}, and -\emph{won} and omit -\emph{aen},
-\emph{waen}, -\emph{ean}, and -\emph{wean} from further discussion
here. presents the type A rimes with dental codas, omitting the vowels
-\emph{ae}- and -\emph{ea}-.

§b. The rime -\emph{wen} occurs only with velar initials (牙 \emph{K}-).
As was the case with the velar finals (cf. § and §), the fact that K- is
compatible with a greater number of rimes suggests that the labial
feature can be credited to the initial rather than the rime. The dental
final rimes do not present strong motivation for labiovelar initials,
but as labiovelar initials are already posited on firm grounds, there is
no reason not to take advantage of them here also. Reformulating with
cognizance of labiovelars results in ; although the elegance of the
distribution is improved, many gaps remain.

\begin{figure}
    \centering
\begin{tabular}{|l|l|l|l|l|l|l|l|}
 \hline
&&
P- &
T- &
Ts- &
K- \\ \hline
{\emph{-an}} &
寒一開 &
{\emph{Pan}} &
{\emph{Tan}} &
{\emph{Tsan}} &
{\emph{Kan}} \\ \hline
{\emph{-wan}} &
桓一合 &
--- &
{\emph{Twan}} &
{\emph{Tswan}} &
{\emph{Kwan}} \\ \hline
{\emph{-en}} &
先四開 &
{\emph{Pen}} &
{\emph{Ten}} &
{\emph{Tsen}} &
{\emph{Ken}} \\ \hline
{\emph{-wen}} &
先四合 &
--- &
--- &
--- &
{\emph{Kwen}} \\ \hline
{\emph{-on}} &
痕一開 &
--- &
--- &
--- &
{\emph{Kon}} \\ \hline
{\emph{-won}} &
魂一合 &
{\emph{Pwon}} &
{\emph{Twon}} &
{\emph{Tswon}} &
{\emph{Kwon}} \\ \hline
\end{tabular}
    \caption{Middle Chinese division-i/iv syllable types in final -n \citep[200]{Baxter2014OldChinese}}
    \label{fig:MChi-div-i-iv-final-n}
\end{figure}

\begin{figure}
    \centering
\begin{tabular}{|l|l|l|l|l|l|l|l|}
 \hline
&
P- &
T- &
Ts- &
K- &
K{\sil{ʷ}}- \\ \hline
{\emph{-an}} &
{\emph{Pan}} &
{\emph{Tan}} &
{\emph{Tsan}} &
{\emph{Kan}} &
\textbf{*K{\sil{ʷ}}an > {\emph{Kwan}}} \\ \hline
{\emph{-wan}} &
--- &
{\emph{Twan}} &
{\emph{Tswan}} &
{\emph{Kwan}} &
\textbf{*K{\sil{ʷ}}wan > {\emph{Kwan}}} \\ \hline
{\emph{-en}} &
{\emph{Pen}} &
{\emph{Ten}} &
{\emph{Tsen}} &
{\emph{Ken}} &
\textbf{*K{\sil{ʷ}}en > {\emph{Kwen}}} \\ \hline
{\emph{-on}} &
--- &
--- &
--- &
{\emph{Kon}} &
\textbf{*K{\sil{ʷ}}on > {\emph{Kwon}}} \\ \hline
{\emph{-won}} &
{\emph{Pwon}} &
{\emph{Twon}} &
{\emph{Tswon}} &
{\emph{Kwon}} &
\textbf{*K{\sil{ʷ}}won > {\emph{Kwon}}} \\ \hline
\end{tabular}
    \caption{Labiovelar initials (division-i/iv in final -n)}
    \label{fig:labiovelar-initials}
\end{figure}

\paragraph{The `rounded vowel' hypothesis:
origin of the remaining -\emph{w}- (合口 \emph{hékŏu}) rimes.} 
\label{par:rounded-vowel}
In the case of type A velars
finals, positing labiovelar initials explained away rimes with medial
-\emph{w}- altogether. However, in the case of dental finals the
labio-velars initials offer no explanation of syllables such as
 \emph{Twan} and \emph{Tswan}. Since medial -\emph{w}- is not required
for velar final syllables, a source for these syllables that makes no
recourse to medial -w- desireable.

Since the Old Chinese rimes *-aŋ, *-oŋ, *-eŋ, *-əŋ, and *-uŋ are set up
for the velar finals the null hypothesis is that only *-an, *-on, *-en,
*-ən, and *-un are to be set up for the dental finals. We may presume
that Middle Chinese -\emph{an}, -\emph{en}, and -\emph{on} {[}ʌn] are
retentions of Old Chinese *an, *en, and *ən respectively. This leaves
Middle Chinese -\emph{wan} and -\emph{won} {[}wʌn] to be explained as
somehow descended from Old Chinese *-on and *-un. Jaxontov proposed that
-\emph{wan} results from the breaking of *-on and -\emph{won} {[}wʌn]
similarly derives from *-un (1960b esp. p. 104, 1970 esp. p. 54, also
see Baxter 1992: 236-240 and \cite[203-207]{Baxter2014OldChinese}). This
internal reconstruction makes the following predictions. (1) Syllables
like \emph{twan} \textless{} *ton will not rhyme with syllables like
\emph{tan} in the 詩経  \emph{Shījīng}. (2) Syllables like \emph{kwan}
will divide into two rhyme classes, those deriving from *kon and those
deriving from *{\sil{kʷ}}an, and members of these two classes with not rhyme
with other. Since these patterns reflect what does not rhyme rather what
does, they cannot easily be illustrated anecdotally. Nonetheless, Baxter
provides detailed evidence that predictions are largely borne out in the
rhyming practices of the 詩經  \emph{Shījīng} (1992: 375-389). Accepting
the rounded vowel hypothesis allowed to be rewritten as .

\begin{figure}
    \begin{adjustbox}{max width=1.3\textwidth,center}
\begin{tabular}{|l|l|l|l|l|l|l|l|}
 \hline
&
P- &
T- &
Ts- &
K- &
Kʷ- \\ \hline
{\emph{-an}} &
{\emph{Pan}} &
{\emph{Tan}} &
{\emph{Tsan}} &
{\emph{Kan}} &
*Kʷan > {\emph{Kwan}} \\ \hline
{\emph{-wan}} < *on &
--- &
\textbf{*Ton > {\emph{Twan}}} &
\textbf{*Tson > {\emph{Tswan}}} &
\textbf{*Kon > {\emph{Kwan}}} &
\textbf{*Kʷon > *Kʷwan > {\emph{Kwan}}} \\ \hline
-en &
Pen &
Ten &
Tsen &
Ken &
*Kʷen > {\emph{Kwen}} \\ \hline
{\emph{-on}} < *ən &
--- &
--- &
--- &
*Kən > {\emph{Kon}} &
*Kʷən > *Kʷon > {\emph{Kwon}} \\ \hline
-won < *un &
\textbf{*Pun > {\emph{Pwon}}} &
\textbf{*Tun > {\emph{Twon}}} &
\textbf{*Tsun > {\emph{Tswon}}} &
\textbf{*Kun > {\emph{Kwon}}} &
\textbf{*Kʷun > *Kʷwon > {\emph{Kwon}}} \\ \hline
\end{tabular}
\end{adjustbox}
    \caption{Breaking of rounded vowels (division-i/iv in final -n)}
    \label{fig:breaking-vowels}
\end{figure}

\paragraph{w-neutralization.} Y. R. Chao noticed that Middle Chinese does not distinguish syllables
like \emph{pan} and syllables like \emph{pwan} (1941; 208; 
\cite[252-256, 567]{Baxter1992}), we can presume that after the change *-on \textgreater{}
-\emph{wan}, syllables like \emph{pwan} merged with \emph{pan}, because
of the labialness of the initial. This hypothesis allows us to plug two
more holes in the distribution of vowels before the dental nasal and
makes the prediction that early rhyming practices will distinguish
\emph{pan} \textless{} *pan from \emph{pan} \textless{} *pwan
\textless{} *pon and \emph{pwon} \textless{} *pun from \emph{pwon}
\textless{} *pən. The rhyming practice of the 詩経  \emph{Shījīng} bears
out this prediction. For example, 奔 \emph{pwon} (33-28a) and 門
\emph{mwon} (33-35a), although they have the same Middle Chinese rime
category, do not rhyme in the 詩経  \emph{Shījīng}, neither does any
other word rhyme with both of them (cf. ).\footnote{A distinction of *p-
  versus *pw-, subsequently lost, is also reconstructed for proto-Tai
  (Pittayaporn 2009: 155) and proto-Tupi (Jensen 1999: 140-141).}
Accepting the w-neutralization hypothesis allowed to be rewritten as .

\begin{figure}[t]
    \begin{adjustbox}{max width=1.3\textwidth,center}
\begin{tabular}{|llll|llll|}
\hline
\multicolumn{4}{|c|}{Words rhyming with 奔 
{\emph{pwon}} (33-28a)} &
\multicolumn{4}{c|}{Words rhyming with 門 mwon (33-35a)} \\ \hline
鶉 &
{\emph{dzywin}} &
34-18j &
(Ode 49.2) &
殷 &
{\emph{'j{\sil{ɨ}}n}} &
33-09a &
(Ode 40.1) \\ 
君 &
{\emph{kjun}} &
34-12a &
(Ode 49.2) &
貧 &
{\emph{bin}} & 
33-30v &
(Ode 40.1) \\
璊 &
{\emph{mwon}} &
24-57f &
(Ode 73.2) &
艱 &
{\emph{kean}} &
33-05c &
(Ode 40.1, 199.1) \\
啍 &
{\emph{thwon}} &
34-18t &
(Ode 73.2) &
雲 &
{\emph{hjun}} &
34-14b &
(Ode 93.1, 261.4) \\
&
&
&
&
存 &
{\emph{dzwon}} &
33-22a &
(Ode 93.1) \\
&
&
&
&
巾 &
{\emph{kin}} &
33-06a &
(Ode 93.1) \\
&
&
&
&
員 &
{\emph{hjun}} &
23-10a &
(Ode 93.1) \\ 
&
&
&
&
云 &
{\emph{hjun}} &
34-14a &
(Ode 199.1) \\ \hline
\end{tabular}
\end{adjustbox}
    \caption{the non-rhyming of 奔 pwon (33-28a) and 門 mwon (33-35a)}
    \label{fig:pwon-and-mwon}
\end{figure}

\begin{figure}[t]
    \centering
\begin{tabular}{|l|l|l|l|l|}
\hline
&
P- &
Tr- &
K- &
Tsy- \\ \hline
-jon &
{\emph{pjon}} &
—  &
{\emph{kjon}} &
— \\ \hline
-jwon &
{\emph{pjwon}} &
— &
{\emph{kjwon}} &
— \\ \hline
-j{\sil{ɨ}}n &
— &
— &
{\emph{kj{\sil{ɨ}}n}} &
— \\ \hline
-jun &
{\emph{pjun}} &
— &
{\emph{kjun}} &
— \\ \hline
-in &
{\emph{pin}} &
{\emph{trin}} &
{\emph{kjin}} &
{\emph{tsyin}} \\ \hline
-jwin &
— &
{\emph{trwin}} &
{\emph{kwin}} &
{\emph{twywin}} \\ \hline
-jen &
{\emph{pj(i)en}} &
{\emph{trjen}} &
{\emph{kj(i)en}} &
{\emph{tsyen}} \\ \hline
-jwen &
— &
{\emph{trjwen}} &
{\emph{kjwen}} &
{\emph{tsjwen}} \\
\hline
\end{tabular}
    \caption{the distribution of Middle Chinese type B rimes before -n}
    \label{fig:distibution-of-MChi-n}
\end{figure}


\paragraph{`a-raising', `acute fronting', and
'{\sil{ɨ}}-raising'.} Before further addressing the
asymmetries in the distribution of type A syllables let us turn our
attention to the asymmetries in the distribution of type B syllables.
presents the distribution of initials with rimes in type B syllables.
The absence of syllables of the shape *pjwin, *p{\sil{jɨ}}n, and *pjwen is no
surprise; as a consequence of w-neutralization (cf. §) these syllables
are not distinguishable from \emph{pin},  \emph{pjun}, and \emph{pjen}.
The more surprising gap is the absence of the rimes -\emph{jon},
-\emph{jwon}, -\emph{{\sil{jɨ}}n}, and -\emph{jun} after acute initials. This
gap suggests that the back and mid vowels have fronted between an acute
initial and a final dental.

§a. Before considering Baxter's formulation of this process of raising,
it is convenient to point out that final *-an does not occur in type B
syllables. Baxter proposes that a fairly recent change, `a-raising'
(1992: 579), changed *-an to -\emph{on} {[}ʌn] in type B syllables.
With this, effectively orthographic, change in mind, the noteworthy
asymmetry of type B syllables becomes the absence of *‑an, *‑wan, *‑ən,
and *‑un after acute initials.

§b. Baxter proposes two separate but related sound changes to account
for this odd distribution: *-an \textgreater{} -en ('acute fronting',
\cite[574-575]{Baxter1992}), which only took place in type B syllables, and
*-ən \textgreater{} -in ('{\sil{ɨ}}-fronting', 
\cite[577]{Baxter1992}), which took
place in both type A and type B syllables, although our motivation for
it so far only applies to type B syllables.

\begin{figure}
    \begin{adjustbox}{max width=1.3\textwidth,center}
\begin{tabular}{|l|lll|ll|ll|ll|lll|}
\hline
&
P- & & &
T- & &
Ts- & &
K- & &
K{\sil{ʷ}}- & & \\ \hline
-an & 
\tikzmark{Pan-b}{\emph{Pan}} &&&
{\emph{Tan}} &&
{\emph{Tsan}} &&
{\emph{Kan}}  &&
*K{\sil{ʷ}}an > & {\emph{Kwan}} & \\ \hline
{\emph{-wan}} < *on &
*\textbf{Pon} > & \textbf{*Pwan}  > &
\tikzmark{Pan-a}\textbf{{\emph{Pan}}} &
*Ton > & {\emph{Twan}} &
*Tson > & {\emph{Tswan}} &
*Kon > & Kwan &
*K{\sil{ʷ}}on > & *K{\sil{ʷ}}wan > & {\emph{Kwan}} \\ \hline
{\emph{-en}} &
{\emph{Pen}} &&&
\tikzmark{Ten-f}{\emph{Ten}} &&
\tikzmark{h}{\emph{Tsen}} &&
{\emph{Ken}} &&
*K{\sil{ʷ}}en > & {\emph{Kwen}} & \\ \hline
-on < *ən &
\textbf{*Pən} > & \textbf{*Pwən} > & \tikzmark{Pwon-c}\textbf{{\emph{Pwon}}}  &
\textbf{*Tən} > & \tikzmark{Ten-e}\textbf{{\emph{Ten}}}
&
\textbf{*Tsən} > & \textbf{{\emph{Tsen}}}\tikzmark{Tsen-g} &
*Kən > & {\emph{Kon}} &
*K{\sil{ʷ}}ən > & *K{\sil{ʷ}}on > & {\emph{Kwon}} \\ \hline
-won < *un &
*Pun > & \tikzmark{Pwon-d}{\emph{Pwon}} &&
*Tun > & {\emph{Twon}} &
*Tsun > & {\emph{Tswon}} &
*Kun > & {\emph{Kwon}} &
*K{\sil{ʷ}}un > & *K{\sil{ʷ}}won > & {\emph{Kwon}} \\ \hline
\end{tabular}
\end{adjustbox}
    \caption{hi > mid (division-i/iv in final -n)}
    \label{fig:hi-mid-div}
      \begin{tikzpicture}[overlay, remember picture, shorten >=6pt, shorten <=6pt, transform canvas={yshift=.2\baselineskip}]
    \draw [-> , line width=1pt] ({pic cs:Pan-a}) [bend right] to ({pic cs:Pan-b});
    \draw [->, line width=1pt] ({pic cs:Pwon-c}) [bend left] to ([yshift=5pt]{pic cs:Pwon-d});
    \draw [->, line width=1pt] ([xshift=-20pt,yshift=3pt]{pic cs:Ten-e}) [bend right] to ([xshift=-20pt]{pic cs:Ten-f});
    \draw [->, line width=1pt] ([xshift=-52pt,yshift=5pt]{pic cs:Tsen-g}) [bend right] to ([xshift=-30pt]{pic cs:h});
  \end{tikzpicture}
\end{figure}

\paragraph{hi \textgreater{} mid.}
Returning to the distribution of type A syllables, there are only two
holes remaining in namely Middle Chinese *Ton and *Tson. This absence
raises the question of how Old Chinese *Tən and *Tsən developed, such
that they did not become Middle Chinese *Ton and *Tson respectively. To
suggest that Old Chinese *Tən developed into  \emph{Twon} or \emph{Twan}
is unmotivated; the syllable *Tən offers no potential origin for Middle
Chinese medial -\emph{w}-. It is much more likely that *Tən developed to
 \emph{Ten} or \emph{Tan}. If we recall the proposal `{\sil{ɨ}}-fronting' of the
previous paragraph (cf. §), there is an explanation available for
-\emph{in} arising from *-ən. Unfortunately, we do not want -\emph{in}
but -\emph{en}. Baxter's solution is to propose a change of *-in to
-\emph{en} in type A syllables only; he dubs this change `hi
\textgreater{} mid' 
\citep[578]{Baxter1992}. The change `hi \textgreater{} mid'
occurred after `{\sil{ɨ}}-fronting', i.e. *-ən \textgreater{} *-in
\textgreater{} -\emph{en} in type A syllables. With this proposal in
place it is possible to explain all of the Middle Chinese type A
syllables with final dentals as originating from a version of Old
Chinese with five vowels (cf. ).

\begin{figure}[t]
    \begin{adjustbox}{max width=1.5\textwidth,center}
\begin{tabular}{|l|llll|lll|ll|llll|l|l|lll}
\hline
&
P- &&&&
Tr- &&&
K- &&
K{\sil{ʷ}}- &&&&
Tsy- \\ \hline
*-en &
{\emph{pj(i)en}} &&&
 & {\tikzmark{trjenj}}{\emph{trjen}} &&&
{\emph{kj(i)en}} &&
{\emph{kjwen}} &&&&
{\emph{tsyen}} \\ \hline
*-an &
\tikzmark{panl}*pan > & {\emph{pjon}} &&&
 *tran >  & {\emph{trjen}}{\tikzmark{trjeni}} & &
*kan > & {\emph{kjon}} &
*k{\sil{ʷ}}an > & {\emph{kjwon}}{\tikzmark{w}} &&&
*tan > {\emph{tsyen}} \\ \hline
*-on &
*pon > & *pwan > & \tikzmark{k}*pan & %{\emph{pjon}} 
&
*tron > & *trwan > & {\emph{trjwen}} &
*kon > & *kwan\tikzmark{s} &
*k{\sil{ʷ}}on > & *k{\sil{ʷ}}wan > & kwan\tikzmark{t} > & {\emph{kjwon}}{\tikzmark{x}} &
{\emph{tsywen}} \\ \hline
*-un &
*pun > & *pwən{\tikzmark{pwənm}}  & 
%> {emph{pjun}} 
& &
*trun > & *trwən > & {\emph{trwin}} &
*kun > & *kwən\tikzmark{u} & 
%{\emph{kjun}} → &
*k{\sil{ʷ}}un > & *k{\sil{ʷ}}wən > & kwən\tikzmark{v} > & {\emph{kjun}}{\tikzmark{z}} &
*tun > *twən > {\emph{twywin}} \\ \hline
*-ən &
*pən > & \tikzmark{pwənn}*pwən > &  {\emph{pjun}} &&
*trən > & {\tikzmark{trinq}}{\emph{trin}} & &
*kən > & {\emph{kj{\sil{ɨ}}n}} &
*k{\sil{ʷ}}ən > & *kwən > & {\emph{kjun}}{\tikzmark{y}} &&
*tən > {\emph{tsyin}}{\tikzmark{tsyina}} \\ \hline
*-in &
{\emph{pin}} &&&&
{\tikzmark{trinr}}{\emph{trin}} &&&
{\emph{kjin}} &&
{\emph{kwin}} &&&&
{\emph{tsyin}}{\tikzmark{tsyinb}} \\ \hline
\end{tabular}
\end{adjustbox}
    \caption{six vowels in type B rimes before -n}
    \label{fig:6-vowel-type-B-n}
\begin{tikzpicture}[overlay, remember picture, 
%shorten >=4pt, shorten <=8pt, 
transform canvas={yshift=-.2\baselineskip}]
    \draw [-> , line width=0.7pt, shorten >=6pt, shorten <=8pt] ([yshift=-2pt,xshift=-30pt]{pic cs:trjeni}) [bend right] to ({pic cs:trjenj});
    \draw [->, line width=0.7pt, shorten >=4pt, shorten <=8pt] ([yshift=3pt]{pic cs:k}) [bend right] to ([yshift=4pt,xshift=2pt]{pic cs:panl});
    \draw [->, line width=0.7pt, shorten >=5pt, shorten <=4pt] ([yshift=10pt,xshift=-29pt]{pic cs:pwənm}) [bend right=140,min distance=1cm] to ([yshift=9pt]{pic cs:pwənn});
    \draw [->, line width=0.7pt] ({pic cs:o}) [bend right] to ({pic cs:p});
    \draw [->, line width=0.7pt, shorten >=6pt, shorten <=8pt] ([xshift=-12pt,yshift=13pt]{pic cs:trinq}) [bend left] to ([xshift=-4pt,yshift=15pt]{pic cs:trinr});
    \draw [->, line width=0.7pt,bend left=110,min distance=2.8cm] ([xshift=-54pt,yshift=8pt]{pic cs:s}) [bend left] to ([xshift=-85pt,yshift=10pt]{pic cs:t});
    \draw [->, line width=0.7pt,bend left=130,min distance=2.6cm] ([xshift=-52pt,yshift=3pt]{pic cs:u}) [bend right] to ([xshift=-75pt,yshift=4pt]{pic cs:v});
    \draw [->, line width=0.7pt,bend left=35,min distance=0.5cm] ([xshift=-85pt,yshift=8pt]{pic cs:x}) [bend right] to ([xshift=-75pt,yshift=6pt]{pic cs:w});
    \draw [->, line width=0.7pt,bend left=40,min distance=0.5cm] ([xshift=-77pt,yshift=3pt]{pic cs:z}) [bend left] to ([xshift=-75pt,yshift=5pt]{pic cs:y});
    \draw [->, line width=0.7pt,bend left=50,min distance=0.5cm] ([xshift=-83pt,yshift=7pt]{pic cs:tsyina}) [bend left] to ([xshift=-75pt,yshift=10pt]{pic cs:tsyinb});
  \end{tikzpicture}
\end{figure}


\begin{figure}[t]
    \begin{adjustbox}{max width=1.5\textwidth,center}
\begin{tabular}{|l|lll|ll|ll|ll|lll|}
\hline
& P- &&&
T- &&
Ts- &&
K- &&
K{\sil{ʷ}}- && \\ \hline
*-an &
{\tikzmark{Pan-2}}{\emph{Pan}} &&&
{\emph{Tan}} &&
{\emph{Tsan}} &&
{\emph{Kan}} &&
*K{\sil{ʷ}}an > & {\tikzmark{Kwan-3}}{\emph{Kwan}} & \\ \hline
*-on &
*Pon > & *Pwan > & {\tikzmark{Pan-1}}{\emph{Pan}} &
*Ton > & {\emph{Twan}} &
*Tson > & {\emph{Tswan}} &
*Kon > & {\tikzmark{Kwan-1}}{\emph{Kwan}} &
*K{\sil{ʷ}}on > & *K{\sil{ʷ}}wan > & {\tikzmark{Kwan-2}}{\emph{Kwan}} \\ \hline
*-in &
*Pin > & {\tikzmark{Pen-1}}{\emph{Pen}} &&
*Tin > & {\tikzmark{Ten-1}}{\emph{Ten}} &
*Tsin > & {\tikzmark{Tsen-1}}{\emph{Tsen}} &
*Kin > & {\tikzmark{Ken-1}}{\emph{Ken}} &
*K{\sil{ʷ}}in > & {\tikzmark{Kwen-1}}{\emph{Kwen}} & \\ \hline
*-en &
{\tikzmark{Pen-2}}{\emph{Pen}} &&&
{\tikzmark{Ten-3}}{\emph{Ten}} &&
{\tikzmark{Tsen-3}}{\emph{Tsen}} &&
{\tikzmark{Ken-2}}{\emph{Ken}} &&
*K{\sil{ʷ}}en > & {\tikzmark{Kwen-2}}{\emph{Kwen}} & \\ \hline
*-ən &
*Pən > & *Pwən > & {\tikzmark{Pwon-1}}{\emph{Pwon}} &
*Ton > & {\tikzmark{Ten-2}}{\emph{Ten}} &
*Tson > & {\tikzmark{Tsen-2}}{\emph{Tsen}} &
*Kən > & {\emph{Kon}} &
*K{\sil{ʷ}}ən > & *K{\sil{ʷ}}on > & {\tikzmark{Kwon-3}}{\emph{Kwon}} \\ \hline
*-un &
*Pun > & {\tikzmark{Pwon-2}}{\emph{Pwon}} &&
*Tun > & {\emph{Twon}} &
*Tsun > & {\emph{Tswon}} &
*Kun > & {\tikzmark{Kwon-1}}{\emph{Kwon}} &
*K{\sil{ʷ}}un > & *K{\sil{ʷ}}won > & {\tikzmark{Kwon-2}}{\emph{Kwon}} \\ \hline
\end{tabular}

\end{adjustbox}
    \caption{six vowels in type A rimes before -n}
    \label{fig:6vowel-in-type-A-n}
          \begin{tikzpicture}[overlay, remember picture, shorten >=6pt, shorten <=6pt, transform canvas={yshift=.2\baselineskip}]
    \draw [-> , line width=1pt] ([xshift=12pt,yshift=2pt]{pic cs:Pan-1}) [bend right] to ([xshift=15pt,yshift=2pt]{pic cs:Pan-2});
    \draw [-> , line width=1pt] ([xshift=12pt,yshift=0pt]{pic cs:Pen-1}) [bend left] to ([xshift=15pt,yshift=2pt]{pic cs:Pen-2});
    \draw [-> , line width=1pt] ([xshift=12pt,yshift=0pt]{pic cs:Pwon-1}) [bend left] to ([xshift=20pt,yshift=2pt]{pic cs:Pwon-2});
    \draw [-> , line width=1pt] ([xshift=12pt,yshift=0pt]{pic cs:Ten-1}) [bend left] to ([xshift=20pt,yshift=2pt]{pic cs:Ten-3});
    \draw [- , line width=1pt] ([xshift=12pt,yshift=0pt]{pic cs:Ten-2}) [bend right] to ([xshift=20pt,yshift=1pt]{pic cs:Ten-3});
   \draw [-> , line width=1pt] ([xshift=12pt,yshift=0pt]{pic cs:Tsen-1}) [bend left] to ([xshift=20pt,yshift=2pt]{pic cs:Tsen-3});
   \draw [- , line width=1pt] ([xshift=12pt,yshift=0pt]{pic cs:Tsen-2}) [bend right] to ([xshift=23pt,yshift=0.5pt]{pic cs:Tsen-3});
    \draw [-> , line width=1pt] ([xshift=12pt,yshift=0pt]{pic cs:Ken-1}) [bend left] to ([xshift=20pt,yshift=2pt]{pic cs:Ken-2});
    \draw [<-> , line width=1pt] ([xshift=30pt,yshift=2pt]{pic cs:Kwen-1}) [bend left=90,min distance=2cm] to ([xshift=30pt,yshift=-2pt]{pic cs:Kwen-2});
    \draw [-> , line width=1pt] ([xshift=12pt,yshift=0pt]{pic cs:Kwon-1}) [bend right=140,min distance=2.4cm] to ([xshift=20pt,yshift=2pt]{pic cs:Kwon-3});
    \draw [-> , line width=1pt] ([xshift=12pt,yshift=0pt]{pic cs:Kwon-2}) [bend right=160,min distance=2.5cm] to ([xshift=20pt,yshift=-2pt]{pic cs:Kwon-3});
    \draw [-> , line width=1pt] ([xshift=12pt,yshift=0pt]{pic cs:Kwan-2}) [bend right=100,min distance=1.5cm] to ([xshift=20pt,yshift=-2pt]{pic cs:Kwan-3});
    \draw [- , line width=1pt] ([xshift=12pt,yshift=0pt]{pic cs:Kwan-1}) [bend left=100,min distance=1.7cm] to ([xshift=21pt,yshift=-2pt]{pic cs:Kwan-3});
\end{tikzpicture}
\end{figure}

\paragraph{Asymmetries in type B rimes with
dental codas.} 
\label{par:Asymmetries-in-type-B-rimes-with-dental-codas}
The distribution of Middle Chinese type A syllables
before velars (cf. §§-) and dentals (cf. §§-) gives rise to a five vowel
theory of Old Chinese phonology, viz. *e, *a, *o, *u, and *ə. Can this
theory, with its attendant hypotheses (labio-velars, cf. §, rounded
vowels, cf. §, w-neutralization, cf. §, a-raising, cf. §a, acute
fronting, cf. §b, {\sil{ɨ}}-fronting, cf. §b, hi \textgreater{} mid, cf. §),
also explain the distribution of Middle Chinese type B syllables, e.g.
before final -\emph{n}? The answer is no. To explain syllables of the
shape \emph{pin} and \emph{trin} one must presume a sixth vowel *-i-
(cf. ). In the type A syllables it is the change `hi \textgreater{} mid'
(cf. §) which obscured *-i- as a separate vowel (cf. ). A reader seeking
assurance that type A *-in and type A *-en rhyme separately may rest
secure in the knowledge that already the Qīng philologists recognized
this distinction, in traditional terminology it is the difference
between the 元 \emph{yuán} and the 真  \emph{zhēn} rhyme groups 
\citep[203]{Baxter2014OldChinese}. In summary, it is judicious to reconstruct Old
Chinese with six nuclear vowels *a, *e, *i, *o, *u, and *ə.

\section{Origins of the tones and final
clusters}

\paragraph{} 
\label{origins-of-the-tones-and-final-clusters}
Old Chinese lacked tones; the
Middle Chinese tones (cf. §) originating from Old Chinese segments. The
distribution of Middle Chinese tones in  \emph{xiéshēng} series and the
 \emph{Shījīng} singles out the departing tone (\emph{qùshēng} 去聲,
-\emph{H}) as potentially secondary (§). Relying on his own pre­vious
work on the origin of tones in Vietnamese, Haudricourt explains that the
departing tone originated from segmental *-s or the final clusters *-ps,
*-ts, and *-ks (§). Consensus holds that the Middle Chinese rising tone
(\emph{shǎngshēng} 上聲, -\emph{X}) originates from *-ʔ, but the
evidence for this proposal is weaker than for the origin of the
departing tone (§).

\paragraph{Distribution of Middle Chinese
tones in  \emph{\textbf{xiéshēng} series and the
 \emph{Shījīng}.}} 
\label{par:distribution-of-MC-tones-in-xs-series-and-the-SJ} 
The phonetic components of Chinese
character generally fall into two natural classes: (1) those appearing
with characters in the level (\emph{píngshēng} 平聲, -Ø) or rising tones
(\emph{shǎngshēng} 上聲, -\emph{X}) in Middle Chinese, and (2) those
appearing in characters in entering (\emph{rùshēng} 入聲, -\emph{p},
-\emph{k}, or -\emph{t}) tone in Middle Chinese. Phonetics of both types
often serve for characters in the departing (\emph{qùshēng} 去聲,
-\emph{H}) tone.

§a. The series built on 反 (24-49) is an example of a series that
combines level, rising, and departing tones.

\begin{description} \item
24-49a 反  \emph{baenX},  \emph{pjonX}

\item
24-49i 飯  \emph{bjonH, bjonX}

\item
24-49j 板  \emph{paenX}

\item
24-49m 扳  \emph{paen}
\end{description}

§b. The series built on 各 (02-01) is an example of a series that
combines entering and departing tones.

\begin{description} \item
02-01a 各 \emph{kak}
\item
02-01f 閣 \emph{kak}
\item
02-01z 格 \emph{kaek}
\item
02-01k' 賂 \emph{luH}
\item
02-01l' 路 \emph{luH}
\end{description}

§c. Duàn Yùcái 段玉裁 (1735-1815), investigating the 詩經  \emph{Shījīng}
rhymes, observed that rhyme sequences in the level, rising, and
departing tones could be isolated, but the rhyming of departing tone
words appeared to be random. He concluded that 古無去聲 ``there was no
departing tone in Old Chinese'' (Sagart 1999b: 96). Thus, for example,
in the following poem from the \emph{Shījīng} (Ode 209) the last
character in each line has a middle Chinese entering tone reading, but
three departing tones interrupt this pattern in an arbitrary way.

\begin{description} \item
執爨踖踖 \emph{tshjek}
\end{description}

\begin{description} \item
為俎孔碩 \emph{dzyek}
\end{description}

\begin{description} \item
或燔或炙 \emph{tsyaeH}
\end{description}

\begin{description} \item
君婦莫莫 \emph{meak}
\end{description}

\begin{description} \item
為豆孔庶 \emph{syoH}
\end{description}

\begin{description} \item
為賓為客 \emph{khaek}
\end{description}

\begin{description} \item
獻酬交錯 \emph{tshak}
\end{description}

\begin{description} \item
禮儀卒度 \emph{duH}
\end{description}

\begin{description} \item
笑語卒獲 \emph{hweak}
\end{description}

\begin{description} \item
神保是格 \emph{kaek}
\end{description}

\begin{description} \item
報以介福 \emph{pjuwk}
\end{description}

萬壽攸酢 \emph{dzak}

\begin{figure}
    \centering
\begin{tabular}{|l|l|l|}
\hline
Mon-Khmer Modern & 
Old Vietnamese &
Vietnamese tone \\ \hline
-s, -h &
-h &
hỏi-ngã \\ \hline
-ʔ &
-ʔ &
sắc-nặng \\ \hline
0 &
0 &
ngang-huyền \\ \hline
\end{tabular}
    \caption{the Mon-Khmer origins of Vietnamese tones}
    \label{fig:MK-origins-Viet-tones}
\end{figure}

\paragraph{Origin of the departing
tone.} Building on his explanation for the origin of
Vietnamese \emph{hỏi-ngã }and \emph{sắc-nặng} tones from laryngeal
endings -h and -ʔ 
\citep[]{Haudricourt1954a}
(Haudricourt 1954a, cf. ), Haudricourt pointed out
that in the oldest layer of Chinese loan-words into Vietnamese, words in
the Chinese departing tone correspond to Vietnamese words in the
 \emph{hỏi-ngã} tone (Haudricourt 1954b)
 \citep[]{Haudricourt1954b}. This suggests that the origin
of the Chinese departing tone is the same as that of the Vietnamese
 \emph{hỏi-ngã} tone. Haudricourt proposed that the departing tone arose
from a suffix -s.

§a. The explanation for departing tone readings appearing in phonetic
series with or rhyming with entering tone readings is the simplification
of final clusters. In the case of velar finals the change in question is
*-ks \textgreater{} *-s \textgreater{} -\emph{H}. The phonetic series
built on 各 (02-01) exhibits this change.

\begin{description} \item
02-01a 各 \emph{kak} \textless{} *{\sil{kˤ}}ak

\item
02-01f 閣 \emph{kak} \textless{} *{\sil{kˤ}}ak

\item
02-01z 格 \emph{kaek} \textless{} *{\sil{kˤ}}rak

\item
02-01k' 賂 \emph{luH} \textless{} *r{\sil{ˤ}}ak-s

\item
02-01l' 路 \emph{luH} \textless{} *Cə.r{\sil{ˤ}}ak-s
\end{description}

A poem from the \emph{Shījīng} (Ode 209) also demonstrates the power of
this hypothesis to improve the interpretation of rhyming in early
poetry.

\begin{description} \item
執爨踖踖  \emph{tshjek} \textless{} *ts{\sil{ʰ}}ak

\item
為俎孔碩  \emph{dzyek} \textless{} *dak

\item
或燔或炙  \emph{tsyaeH} \textless{} *tak-s

\item
君婦莫莫  \emph{meak} \textless{} *mr{\sil{ˤ}}ak

\item
為豆孔庶  \emph{syoH} \textless{} *stak-s

\item
為賓為客  \emph{khaek} \textless{} *{\sil{kʰrˤ}}ak

\item
獻酬交錯  \emph{tshak} \textless{} *ts{\sil{ʰ}}{\sil{ˤ}}ak

\item
禮儀卒度  \emph{duH} \textless{} *l{\sil{ˤ}}ak-s

\item
笑語卒獲  \emph{hweak} \textless{} *mq{\sil{ʷˤ}}rak

\item
神保是格  \emph{kaek} \textless{} *kr{\sil{ˤ}}ak

\item
報以介福  \emph{pjuwk} \textless{} *pək

\item
萬壽攸酢  \emph{dzak} \textless{} *dz{\sil{ˤ}}ak
\end{description}

§b. In the case of dental finals, the final does not disappear entirely,
but leaves a trace into Middle Chinese as final -\emph{j. }The change in
question is *-ts \textgreater{} *-js \textgreater{} -\emph{jH}, as shown
in the phonetic series built on 氒 (22-01).

\begin{description} \item
22-01h 括  \emph{khwat} \textless{} *{\sil{kʷˤ}}at

\item
22-01m 活  \emph{hwat}\textless{} *g{\sil{ʷˤ}}at

\item
22-01o 話 \emph{hwaejH} \textless{} *g{\sil{ʷ}}r{\sil{ˤ}}at-s
\end{description}

§c. In the case of labial finals the early change *-ps \textgreater{}
*-ts lead to a merger of original *‑ps and *-ts. After this merger *-ts
developed as immediately described (i.e. *-ts \textgreater{} *-js
\textgreater{} -\emph{jH}). The early date of the change*‑ps and *-ts
means that readings with final *-ts were available within a phonetic
series when characters were still being coined. As a result phonetic
series with final *-p sometimes include readings with final *-t, as is
the case with the series built on 入 (37-16a).

\begin{description} \item
37-16a 入  \emph{nyip} \textless{} *nup
\end{description}

\begin{description} \item
37-16e 内  \emph{nwojH} \textless{} *n{\sil{ˤ}}uts \textless{} *n{\sil{ˤ}}ups
\end{description}

\begin{description} \item
37-16j 訥 \emph{nwot} \textless{} *n{\sil{ˤ}}ut
\end{description}

§d. Early Chinese transcriptions of foreign words sometimes show the
departing tone cor­re­spond­ing to foreign -\emph{s} ().\footnote{The
  examples are from Pulleyblank (1962: 217-219); I provide the Han
  Chinese reconstructions for reference, based on Schuessler 2009. These
  first seven examples Pulleyblank collected from an article by Bailey
  (1946). The following three examples are from Lokakṣema's (支婁迦讖
  circa 147-189 CE) translation of the \emph{Prajñāpāramitā Sūtra}
  (\emph{Dào Xíng Bānruò Jīng} 道行般若經, T. 224, cited after
  Pulleyblank). The remaining examples come from early historical
  sources. The example of 劫貝 for Skt.  \emph{kārpāsa} `cotton' is from
  Schuessler (2009: 22). In the case of the Skt. examples one must keep
  in mind that the final -a in Sanskrit is likely to have already become
  mute before these words were brought to China. } In addition, some
early Chinese loanwords into Korean have -\emph{s} for Middle Chinese
departing tone (Zhèngzhāng 2000: 63, cf. ).\footnote{Schuessler objects
  to the blanket treatment of 去聲  \emph{qùshēng} as -s on the basis of
  foreign transcriptions, pointing out that all of the examples in are
  cases of -\emph{jH} \textless{} *-ts (2009: 23). One may note that the
  Sino-Korean evidence is also of the same type. He finds use of
  characters with non-dental final 去聲  \emph{qùshēng} readings, in the
  rare cases they are used in foreign transcriptions, correspond to -h,
  -x, or an omission of a final syllable in the foreign original (2009:
  23). Schuessler consequently reconstructs *-(t)s and *h as two
  distinct sources of 去聲 \emph{qùshēng}, albeit in complementary
  distribution. For velar final words, Schuessler's *-kh is implausible,
  since it posits the unlikely series of changes -ks \textgreater{} -{\sil{kʰ}}
  \textgreater{} -h rather than the more straightforward -ks
  \textgreater{} -s \textgreater{} -h. incorporates corrections
  suggested by Sven Osterkamp (\emph{per litteras} 12 February 2017).}

\begin{figure}[t]
    \centering
\begin{tabular}{|llll|}
\hline
Character &
Middle Chinese &
Hàn Chinese &
Foreign language transcribed \\ \hline
波羅奈 &
pa la najH &
{\sil{pɑi lɑi nɑs}} & 
Skt. 
{\emph{vārāṇasī}} \\
三昧 &
sam mwojH &
{\sil{sɑm məs}} & 
Skt. {\emph{samādhi}} (Skt. dh- > Prakrit z-) \\
提謂 &
tej hjwɨjH &
{\sil{de wus}} &
Skt. {\emph{trapuṣa}}, Khot. {\emph{tträväysa}} \\
忉利 &
taw lijH &
{\sil{tɑu lis}} &
Skt. {\emph{trāyastriṃśa}}, Khot. {\emph{ttāvatrīśa}} \\
阿魏 &
'a ngjwɨjH &
{\sil{ʔɑ ŋuih}} & \makecell{
Khot. {\emph{aṃguṣḍä}}, Tokh. B: {\emph{ankwaṣ}}, \\ Uighur {\emph{'nk'pwš}} `asafoetida'} \\
舍衞 &
syaeX hjwejH &
{\sil{śaʔ was}} &
Skt. {\emph{śrāvastī}} \\
迦維羅衛 &
kae ywij la hjwejH &
{\sil{kai wi lɑi was}} & 
Skt. {\emph{kapilavastu}} \\
阿會亘 & 
'a hwajH hwan &
{\sil{ʔɑ ɣuɑs huar}} &
Skt. {\emph{abhāsvara}} \\
阿迦貮吒 &
'a kae nyijH traeH & 
{\sil{ʔɑ ka ńis ṭah}} &
Skt. {\emph{akaniṣṭha}} \\
首陀衛 &
syuwX da hjwejH &
{\sil{śuʔ dɑi was}} &
Skt. {\emph{śuddhāvāsa}} \\
對馬 &
twojH maeX &
{\sil{tuəs maʔ}} &
Jp. Tusima (modern Tsushima) \\
都賴 &
tu lajH &
{\sil{tɑ lɑs}} & 
Talas \\
罽賓 &
kjejH pjien &
{\sil{kɨɑs pin}} & 
Kashmir \\
蒲纇 &
bu lwojH &
{\sil{bɑ luis}} &
Mong. {\emph{bars}} `tiger'  \\
劫貝 &
kjaep pajH &
{\sil{kɨɑp pɑs}} &
Skt. {\emph{kārpāsa}} `cotton' \\ \hline
\end{tabular}
    \caption{the 去聲 qùshēng tone transcribing foreign -s}
    \label{fig:qusheng-as-s}
\end{figure}

\begin{figure}[t]
    \centering
\begin{tabular}{|lllll|}
\hline
Character &
Middle Chinese &
Korean &
& \\ \hline
味 &
mj{\sil{ɨ}}jH &
{\kor{맛}}
 & 
mas &
'taste' \\
器 &
khijH &
{\kor{그릇}} &
kulus &
`implement, utensil' \\
篦 &
pjijH & 
{\kor{빗}} & 
pis &
`comb’ \\
芥 &
kaejH &
{\kor{갓}} &
kas &
`mustard plant’ \\
蓋 &
kajH &
{\kor{갓}} &
kas &
‘cover, bamboo hat’ \\
製 &
tsyejH &
{\kor{짓}} &
cis- &
`to make, construct'  \\ \hline
\end{tabular}
    \caption{the 去聲 qùshēng tone as Korean -s}
    \label{fig:qusheng-Korean}
\end{figure}

\paragraph{Origin of the rising
tone.} 
\label{par:origin-of-the-rising-tone}
\citet[]{Mei1970}
%Mei (1970) 
proposes that the Chinese rising
tone developed out of an earlier glottal stop ending, like the
Vietnamese \emph{sắc-nặng} tone in Haudricourt's theory. His proposal
helps account for the correspondence between the Chinese rising tone and
the Vietnamese \emph{sắc-nặng} in the early layer of Chinese loan-words
into Vietnamese. In further support of his proposal, Mei points to
Chinese dialects which pronounce a glottal stop in some words written
with characters that have rising tone readings in Middle Chinese (1970:
89). He also notes a comment of 義淨 Yìjìng (635--713 CE) in his
南海寄歸內法 \emph{Nánhǎi jìguī nèifǎ}, which explicitly says that
characters used to transliterate short vowels in Sanskrit 皆須上聲讀之
``should all be read in the rising tone'' (Mei 1970: 90), regardless of
a character's normal reading.

This evidence that the rising tone was originally a glottal stop is not
nearly as convincing as the evidence that the departing tone derives
from final *-s. Nevertheless, I represent the rising tone as -ʔ to
accord with 
\citet[195]{Baxter2014OldChinese}.

\section{Finals of Old Chinese}

\paragraph{}
\label{par:finals-of-old-chinese}
The forgoing analysis of tonogenesis makes the tacit assumption that,
when there is no reason to do otherwise the finals -\emph{k}, -\emph{p},
-\emph{t}, -\emph{j}, -\emph{w}, -\emph{ng} (ŋ), -\emph{m}, -\emph{n}
are projected directly backwards from Middle Chinese to Old Chinese. If
Middle Chinese and Old Chinese displayed the same set of finals,
according to the \emph{xiéshēng} hypothesis all characters in the same
 \emph{xiéshēng} series should show Middle Chinese pronunciations with
the same final. In general, this situation prevails, but two interacting
exceptions require comment. In one case final *-j is present in Old
Chinese, even when lacking in Middle Chinese (cf. §). In the other case
Old Chinese has *-r where Middle Chinese has -\emph{n} or -\emph{j} (or
-Ø \textless{} *-j) (cf. §).

\paragraph{-*aj monophthongization.}
\citet[269]{Baxter2014OldChinese}, following 
\citet[292, 413, 570-571]{Baxter1992}
% Baxter (1992: 292, 413, 570-571)
posit *-aj as the Old Chinese origin of Middle Chinese
-\emph{a}.\footnote{OChi. *-a becomes MChi. -\emph{u} in type A
  syllables and -\emph{o} [ʌ] in type B syllables, cf. \cite[572, 575]{Baxter1992}.}
   The motivation for `-aj monophthongization' is the presence
of final glides in early Chinese loans into other languages and in the
relevant cognates of Mǐn dialects of Chinese. The following examples
give relevant Mǐn evidence from the Foochow dialect.

\begin{description} \item
舵  \emph{daX} (18-09-), Foochow /tuai 6/
\item
磨 \emph{ma} (18-18f), Foochow /muai 2/
\item
我 \emph{ngaX} (18-05a), Foochow /ŋuai 3/
\item
破 \emph{phaH} (18-16o), Foochow /p{\sil{ʰ}}uai 5/
\item
跛 \emph{paX} (18-16m), Foochow /pai 3/
\item
簸 \emph{paH} (18-16n), Foochow /puai 5/
\end{description}

This suggestion means that a final *-j must be added in the appropriate
words in Old Chinese.

\paragraph{Final *-r.} 
\label{par:final-r}
Starostin motivated final *-r in Old Chinese, primarily on the basis
of phonetic series that mix reading that have both final -\emph{an} and
(*aj \textgreater{}) -\emph{a} in Middle Chinese (S. Starostin 1989:
399-401). The series built on 釆 (24-54) is a case in point.

\begin{description} \item
24-54b 番 \emph{pa} \textless{} *p{\sil{ˤ}}aj \textless{} *p{\sil{ˤ}}ar
\item
24-54d 幡 \emph{phjon}
\item
24-54g 轓 \emph{pjon}
\item
24-54n 潘 \emph{phan}
\item
24-54o 蟠 \emph{ban}
\item
24-54r 皤 \emph{ba} \textless{} *b{\sil{ˤ}}aj \textless{} *b{\sil{ˤ}}ar
\end{description}

Baxter \& Sagart provide extensive evidence to argue that *-r
\textgreater{} *-n is the mainstream development whereas *-r
\textgreater{} *-j is typical of an eastern dialect 
\citeyearpar[252-268]{Baxter2014OldChinese}.
Although the overall thrust of Baxter \& Sagart's discussion strongly
suggests *-r as part of the Old Chinese syllable canon, because they
approach the evidence unsystematically it is difficult to be confident
of *-r in any one of their reconstructions (cf. Hill 2017 `Chi. *-r').
Nonetheless, in the absence of a dedicated study of Old Chinese *-r, I
have no choice but to defer to their judgements. The reader should
exercise due caution wherever my argument relies on Old Chinese *-r
(i.e. §§, -, and ).

\section{How to reconstruct a word in Old
Chinese}

\paragraph{}
\label{par:how-to-reconstruct-a-word-in-old-chinese}
Baxter \& Sagart do not present their system in a way that
facilitates others reconstructing Old Chinese words in this system.
However, because their assumptions and methods are largely explicit, it
is possible to distill from their discussion a sort of `recipe' for
reconstructing an Old Chinese word step by step. It is most practical to
approach the reconstruction of the rime (cf. §a) and the onset (cf. §b)
separately.

§a. The rime of an Old Chinese word can be determined with the following
three steps:

\begin{enumerate}
\item  Begin with the Middle Chinese reading.

\item Determine whether the rime has multiple origins.

\item Use  \emph{Shījīng} rhymes or  \emph{xiéshēng} evidence to identify a
unique origin of the rime.
\end{enumerate}

In the first step one identifies the Middle Chinese reading of the word
of interest. In some cases, the Middle Chinese rime unambiguously
descends from only one Old Chinese origin, e.g. the Middle Chinese rime
category 模 -\emph{u} has only Old Chinese vowel *-a as its origin, and
thus the Old Chinese ancestor of 堵 \emph{tuX} must have had the rime
*-aʔ. However, in most cases the origins of a Middle Chinese rime
category is ambiguous. For example, the Middle Chinese rime 麻
-\emph{ae} merges Old Chinese *-ra and *-raj. Therefore, 馬 \emph{mæX}
can \emph{a priori} descend from either *mr{\sil{ˤ}}aʔ or *mr{\sil{ˤ}}ajʔ. Thus, in the
second step one must identify whether the Middle Chinese rime has an
unambiguous single origin in Old Chinese, or whether it has multiple Old
Chinese origins. A list of all Middle Chinese rime categories and their
respective potential origins in Old Chinese appears at § . In the third
step one must decide among the potential ancestral rimes using the
rhymes of the \emph{Shījīng} or  \emph{xiéshēng} evidence. The word 馬
\emph{mæX} `horse' rhymes 21 times in the \emph{Shījīng;} in many cases
the word it rhymes with unambiguously has the rime *-a, e.g. 組
 \emph{tsuX} \textless{} *ts{\sil{ˤ}}aʔ (Ode 78.1). Therefore, we must
reconstruct 馬 *mr{\sil{ˤ}}aʔ. In contrast, 麻  \emph{mae} `hemp', which also has
a Middle Chinese reading that permits the Old Chinese rimes *-ra and
*‑raj, rhymes three times in the \emph{Shījīng}. Two of its rhyme words
are unambiguous, viz. 娑 \emph{sa} \textless{} *s{\sil{ˤ}}aj (137.2) and 歌
 \emph{ka} \textless{} *{\sil{kˤ}}aj (Ode 139.1). Baxter provides a finding list
of all rhyme words in the \emph{Shījīng } according to their
 \emph{pīnyīn} transliteration (1992: 745-812); he also gives the rhyme
words, with their Middle Chinese transcriptions and Old Chinese
reconstructions, in the order that they appear in the \emph{Shījīng
}(1992: 583-743).\footnote{In using Baxter (1992) one must bear in mind
  that he does not reconstruct Old Chinese *-r. }

Only a tiny minority of words appear as rhymes in the \emph{Shījīng}.
For words not appearing in the \emph{Shījīng} only  \emph{xiéshēng}
evidence permits the disambiguation of competing Old Chinese
re­con­struc­tions. 
\citet[]{Schuessler2009minimal}
%Schuessler (2009) 
offers the most convenient
presentation of  \emph{xiéshēng} evidence. One must however bear in mind
that his decisions of how to assemble  \emph{xiéshēng} series are his own
and do not necessarily concur with those of Baxter \& Sagart. Purely by
way of illustration, we consider the disambiguation of the Old Chinese
rimes of 馬 \emph{mæX} `horse' and 麻 \emph{mae} `hemp' using
 \emph{xiéshēng} evidence. Schuessler places 麻  \emph{mae} in series
18-18, in which one also finds 摩 \emph{ma} descending unambiguously
from *-aj 
\citeyearpar[217-218]{Schuessler2009minimal}. Schuessler places 馬  \emph{maeX} in series
01-73, in which all readings have Middle Chinese -\emph{ae}- of
ambiguous origin (2009: 63). Nonetheless, Schuessler reconstructs *-a
rather than *-aj, presumably on the tacit basis of the rhyme evidence
presented immediately above.

§b. The initial of an Old Chinese word can be determined, in so far as
is possible, with the following four steps:

\begin{enumerate}
\item 1. Begin with the Middle Chinese initial.
\item 2. Mechanically rewrite Middle Chinese notation in the corresponding Old
Chinese notation.
\item 3. Consult  \emph{xiéshēng} evidence to distinguish laterals, voiceless
resonants, and uvulars.
\item 4. Consult cognate Mǐn dialect forms and borrowings into other
languages.
\end{enumerate}

In the first step one identifies the Middle Chinese reading of the word
of interest.\footnote{In Baxter \& Sagart's system there are no cases in
  which the Middle Chinese initial unambiguously descends from only one
  Old Chinese origin, since it is always possible that there was a loose
  pre-initial which simply dropped. Nonetheless, some Middle Chinese
  initials have more origins in Old Chinese than others do. For example,
  leaving aside loose pre-initials 端 \emph{t}- has only the two origins
  *t- and *C.t, but 定 \emph{d}- derives from *d-, *lˁ-, *N.t- *m-t,
  etc. }

In the second step, a number of changes can be made that amount to
rewriting the Middle Chinese form in the notation of Old Chinese: Middle
Chinese 匣 \emph{h}- is rewritten as Old Chinese *g- (cf. §); 來
\emph{l}- is rewritten as *r- (cf. §); palatal initials and retroflex
initials are mechanically replaced with dentals (i.e. 章 \emph{tsy}-
\textless{} *t-, etc. cf. §, and 知 \emph{tr}- {[}ʈ] \textless{} *tr-
etc. cf. §); pharyngealization (ˁ) is added to a syllable that lack
-\emph{j- }other than those with the main vowel -i-; -\emph{j}- is
removed from a syllable that has it; the velars and laryngeals initials
of words with the relevant Middle Chinese rimes (viz. 唐阳 ‑ \emph{wang}
/ ‑ \emph{jwang} , 登 ‑ \emph{wong}, 庚二 -\emph{waeng}, 耕 -\emph{weang},
and 先 -\emph{wen}) may be rewritten as labiovelar and labiolaryngeals
(cf. §); initial 云 \emph{hj}- is rewritten as *{\sil{ɢʷ}}- (cf. §b).

In the third step  \emph{xiéshēng} evidence permits certain mergers to be
undone:

\begin{description} \item
幫~ \emph{p}-\emph{ }originates from  \emph{*}p.k-\emph{ }if there are
connections to velar stops and otherwise it continues inherited *p- (cf
§a);
\item
滂~ \emph{ph}-\emph{ }originates from *p.{\sil{qʰ}}{\sil{ˤ}}- if there are connections to
曉~ \emph{x}- (cf. §b), from *p.r̥- if there are connections to
徹~ \emph{trh}- (cf §c), and otherwise it continues inherited *p{\sil{ʰ}}-;
\item
透~ \emph{th}- comes from *{\sil{n̥ˤ}}- if it has con­nec­tions to 泥~ \emph{n}-
(cf. §c), *r̥{\sil{ˤ}}- if it has connections to 來~ \emph{l}- (cf.~§e), or *l̥{\sil{ˤ}}-
if it has con­nec­tions to 以~ \emph{y}- or 書~ \emph{sy}- (cf. §d);
\item
定~ \emph{d}- can be re­con­structed as *lˁ- if it has \emph{xiéshēng} con­nec­tions to initial 以~ \emph{y}- or 書~ \emph{sy}- (cf. §)
otherwise it continues inherited *d-;
\item
知 \emph{tr}- derives from *t.kr- if there are connections to velars
(cf. §a), otherwise it continues *tr-;
\item
徹  \emph{trh}- may descend from *t.{\sil{kʰ}}r- if there are connections to
velars (cf. §c), from *r̥- if it has connections to 來~ \emph{l}- (cf.
§c), and otherwise it continues inherited *t{\sil{ʰ}}r-;
\item
\textbf{清~} \emph{tsh}- originates from *s.{\sil{n̥ˤ}}- in
series where it has con­nec­tions to 泥~ \emph{n}- and *s.l̥- if it has
con­nec­tions to 定~ \emph{d}-, 透~ \emph{th}-, and 以~ \emph{y}- (cf.~§);
\item
從  \emph{dz}- comes from *s.b- if there are  \emph{xiéshēng }connections
to labials (cf.~§), otherwise it continues inherited *dz-;
\item
心~ \emph{s}- derives from *s.m({\sil{ˤ}})-, *s.n({\sil{ˤ}})-, *s.ŋ({\sil{ˤ}})-, *s.l{\sil{ˤ}}-, or
*s.q{\sil{ʷ}}({\sil{ˤ}})- if it has \emph{xiéshēng }connections respectively with 明
\emph{m}- (cf.~§a), 泥 \emph{n}- (cf.~§b), 疑  \emph{ng}- (cf.~§c),
定~ \emph{d}-, 透~ \emph{th}-, and 以~ \emph{y}- (cf.~§d), or 云 \emph{hj}-
and 曉 \emph{x}- (cf.~§e);
\item
邪~ \emph{z-} originates from *s.ɢ- in series that
contain characters with initial 以  \emph{y}- (cf.~§§, b), *sə.l- if it
has con­nec­tions to 定~ \emph{d}-, 透~ \emph{th}-, and 以~ \emph{y}- (cf.
§), and otherwise *s.d- (or *s-m-t-) (cf. § and §).
\item
莊\textbf{~} \emph{tsr}- originates from *s.q{\sil{ˤ}}r- in uvular
series (§), otherwise it continues inherited *tsr-;
\item
初~ \emph{tsrh}- originates from *s.r̥- in type B syllables if there is a
con­nec­tion to 生~ \emph{sr}- and *s.t{\sil{ʰ}}({\sil{ˤ}})r- if there is a connection to
palatals, themselves descending from dentals in type B syllables
(cf.~§);
\item
生~ \emph{sr}- can descend from *s.r- if it has \emph{xiéshēng
}con­nec­tions to 來~ \emph{l}- or *s.ŋr- if it has con­nec­tions to
疑~ \emph{ng}- (cf. ~§);
\item
俟\textbf{~} \emph{zr}- originates from *s.ɢr- in uvular series,
i.e. those mixing 見~ \emph{k}-, 群~g-, 匣~ \emph{h}-, 曉~ \emph{x}-, and
以\textbf{~} \emph{y}- (cf.~§), however their example is only of
contact with 匣 \emph{h}- in a series (04-30, based on 㠯 \emph{yiX})
that general appears to be a lateral series; 
\citet[]{Baxter2014OldChinese}
do not appear to propose other origins for 俟\textbf{~} \emph{zr}-.
\item
昌 \emph{tsyh}- may derive from *t.{\sil{kʰ}}- in type B syllables if there are
connections to velar initials (cf.~§b), otherwise it derives from *t-;
\item
書~ \emph{sy}- can be re­con­structed as *n̥- if it has con­nec­tions to
泥~ \emph{n}- (cf.~§c), *l̥- if it has con­nec­tions to 定~ \emph{d}-,
透~ \emph{th}-, or 以~ \emph{y}- (cf.~§d), or *s.t-, *s.t{\sil{ʰ}}-, *s.k-, and
*s.{\sil{kʰ}}- (the latter two only before front vowels) in type B syllables if
there are connections to dentals or velars respectively (cf.~§);
\item
船\textbf{~} \emph{zy}-\emph{ }originates from \emph{*s.g-
}before front vowels in type B syllables if there are con­nec­tions to
velars (cf.~§);
\item
以~ \emph{y}- can be re­con­struct­ed as *l- if it has con­nec­tions to
定~ \emph{d}-, 透~ \emph{th}- or 書~ \emph{sy}- (cf.~§), otherwise it
descends from *{\sil{ɢ}}- (cf.~§b);
\item
見  \emph{k}- originates from *k.l{\sil{ˤ}}- in type A syllables if the word has
 \emph{xiéshēng} connections to 定~ \emph{d}-, 透~ \emph{th}-, and
以~ \emph{y}- (§a), from {\sil{*k.ʔ-}} if the connections are with 影 `- (cf.
§c), from *C.q{\sil{ˤ}}- if the connections are with 曉  \emph{x}-, 影 `-, or 以
 \emph{y}- (cf. §), and otherwise it continues inherited *k-;
\item
溪~ \emph{kh}- originates from *k.r̥{\sil{ˤ}}-in type A syllables if there are
connections to to 來~ \emph{l}- (b), from *C.{\sil{qʰ}}({\sil{ˤ}})- if it has \emph{
}connections to initials 曉 \emph{x}-, 影~'- {[}ʔ], or 以  \emph{y}-
(cf. §), and otherwise continues inherited *{\sil{kʰ}};
\item
疑 \emph{ng}- can be reconstructed as *m-{\sil{qʰ}}-, *m-ɢ-, *N-{\sil{qʰ}}-, or *N-ɢ- if
the word has \emph{xiéshēng }connections to initials 見 \emph{k}-, 影
\emph{'}-, 以  \emph{y}- (cf. §a);
\item
Middle Chinese initial glottal stop (影 `-) can be re­con­structed as
*q- if it has connections to 見 \emph{k}-, 群 g-, 匣 \emph{h}-, 曉
 \emph{x}-, or 以  \emph{y}- (cf. §), otherwise it continues inherited
{\sil{*ʔ-}}.
\item
曉 \emph{x}(\emph{j})- can be reconstructed as *ŋ̊({\sil{ˤ}})- if it has
connections to 疑~ \emph{ng}- (cf~ §a) or *m̥{\sil{ˤ}}- if it has con­nec­tions to
明~ \emph{m}- (cf. §b), other­wise it descends from {\sil{*qʰ}}- (cf.~§ and §a);
\item
匣 \emph{h}- (i.e. *g- in type A syllables following the mechanical
rewriting of step two) can be re­con­structed as *{\sil{ɢ}}ˁ- if it has
connections to 影 `- {[}ʔ] or 以  \emph{y}- (cf.~§a), otherwise it
descends from *g- (cf.~§).
\end{description}

In principle the fourth step is to consider the evidence of Mǐn,
Hmong-Mien, Vietic languages, etc. Unfortunately, this step is currently
impractical. Schuessler (2009) does not include such information and
\citet[]{Baxter2014OldChinese}
do not present such information in a
comprehensive, lexicographically convenient format. Until the requisite
information is collected in a convenient format, if one wants to
incorporate these sources of evidence the only option is to credulously
use the reconstructions that 
\citet[]{Baxter2014OldChinese}
offer.

\section{From Old Chinese to
Trans-Himalayan}\label{from-old-chinese-to-trans-himalayan}

\paragraph{}Now that the lengthy presentation of Old Chinese and its
reconstruction is in place, attention turns to innovations that Old
Chinese reflects vis-à-vis the Trans-Himalayan proto-language. Because
of the ambiguous and indirect nature of information about the
pronunciation of Old Chinese it is not easy to draw a sharp line between
Old Chinese and forms of speech ancestral to it. The most clear cut
methodological approach is to conceptualize as Old Chinese only as much
as can be known without recourse to comparative data. The application of
this methodological principle means that some of the mergers which
Baxter \& Sagart treat as taking place within the history of Chinese are
presented here (cf. §§-).

All of the changes from Trans-Himalayan to Chinese are mergers and these
changes do not interact with one another. As a consequence it is not
possible to resolve the relative chronology of these changes. The
approach here it to first present those mergers that Baxter \& Sagart
treat as Chinese internal and then turn to changes proposed entirely on
the basis of comparison to Tibetan and Burmese.

\paragraph{Merger of velars and dentals after
-i-.} In Old Chinese it is difficult to distinguish dentals
and velars after the vowel -i- (Baxter 1992: 435-437). Because Burmese
merges dentals and velars after -i-, it is of no help in resolving the
Chinese ambiguity (cf. §).

§a. Examples to be reconstructed *iŋ:

\begin{description}
\item
Chi. 年 \emph{nen} \textless{} *C.n{\sil{ˤ}}i{[}ŋ] (32-28a) `harvest, year',
Tib. {\tib{ན་ནིང་}} \emph{na-niṅ }'last year'
\item
Chi. 薪 \emph{sin} \textless{} *si[n] (32-33n) `firewood', Tib. {\tib{ཤིང་}} \emph{śiṅ} `wood'

\item
Chi. 莘 \emph{srin} \textless{} *sri[n] (32-33h) `long, drawn-out',
Tib. {\tib{རིང་}} \emph{riṅ} `long'

\item
Chi. 仁  \emph{nyin} \textless{} *ni{[}ŋ] (32-28f) `kindness', Tib.{\tib{སྙིང་}} \emph{sñiṅ} `heart'
\end{description}

§b. Examples to be reconstructed *-in:

\begin{description} \item
Chi. 憐 \emph{len} \textless{} *r{\sil{ˤ}}iŋ (32-26l) `love, pity', Tib. {\tib{དྲིན་
}} \emph{drin }'kindness'
\end{description}

\begin{description} \item
Chi. 辛 \emph{sin} \textless{} *sin (32-33a) `pungent, painful', Tib. {\tib{མཆིན་}} \emph{mčhin} `liver'
\end{description}

§c. Examples to be reconstructed *ik

\begin{description} \item
Chi. 縊 \emph{'ejH} \textless{} {\sil{*qˤ}}iks (08-05g) `strangle', Tib. {\tib{འཁྱིག་}} \emph{ḫkhyig}

\item
Chi. 節 \emph{tset} \textless{} *ts{\sil{ˤ}}ik (29-30e) `joint of bamboo', Tib. {\tib{ཚིགས་}} \emph{tshigs} `joint'

\item
Chi. 蝨 \emph{srit} \textless{} *sri{[}k] (29-35a) `louse', Tib. {\tib{ཤིག་}} \emph{śig}
\end{description}

\paragraph{Labial neutralization}. Because of
an early neutralization of labial and non-labial vowels before labial
consonants in Old Chinese, it is difficult to distinguish among *-up,
*-op, *-əp or among *-um, *-om, and *-əm (Baxter 1992: 541-553, 
\cite[306]{Baxter2014OldChinese}). In principle, comparative evidence is available to
support the reconstruction of one vowel or another.

\begin{description} \item
Chi. 入 \emph{nyip} \textless{} *nup (37-16a) `enter', Tib. {\tib{ནུབ་}} \emph{nub} `sink, set'

\item
Chi. 立 \emph{lip} \textless{} *k.rəp (37-15a) `stand (v.)', Tib. {\tib{འཁྲབ་
}} \emph{ḫkhrab} \textless{} *ḫkrəp `strike, stamp, tread heavily', OBur. {\burm{ရျပ်}} \emph{ryap} `stand, stop, halt'

\item
Chi. 三 \emph{sam} \textless{} *srum (38-30a) `three', Tib. {\tib{གསུམ་}} \emph{gsum}, Bur. {\burm{သုံး}} \emph{suṃḥ}

\item
Chi. 含 \emph{hom} \textless{} *Cə-m-{\sil{kˤ}}əm (38-03l') `hold in the mouth',
Tib. {\tib{འགམ}}་ \emph{ḫgam} \textless{} *ḫgəm `put in the mouth'

\item
Chi. 尋 \emph{zim} \textless{}*sə-l{[}ə]m (38-17) `warm up (food)',
Bur. {\burm{လုံ}} \emph{luṃ} `warm'
\end{description}

\paragraph{Inability to distinguish *-u- and
in -ə- before acutes.} In many cases Chinese
has the vowel *-ə where one expects *-u-. Because in Old
Chinese it is difficult to distinguish *‑ən, *‑ər, and *‑əj
from *‑un, *‑ur, and *‑uj (Baxter 1992: 427-428, 550; Hill 2012 `six
vowels' : 18-21), it is likely that this complication is due to an
incorrect reconstruction of the vowels in Old Chinese. 

\begin{description} \item
Chi. 運 \emph{hjunH} \textless{} *{\sil{[ɢ]ʷ}}ər-s \textless{}
*{\sil{[ɢ]ʷ}}ərl-s (cf. §) (23-13d) `move', Tib. {\tib{འགུལ་}} \emph{ḫgul}
\textless{} *ḫɢul (Peiros \& Starostin's law, cf. §) \textless{} *ḫɢurl
(cf. §) `move'
\item
Chi. 眉 \emph{mij} \textless{} *mrər \textless{} *mrərl (cf. §) (27-14a)
'eyebrow', WBur. {\burm{မွေး}} \emph{mweḥ} \textless{} *muyḥ \textless{} *murlḥ
(cf. §) `body hair'
\item
Chi. 根 \emph{kon} \textless{} *[k]{\sil{ˤ}}ə[n] \textless{}
*[k]{\sil{ˤ}}ərl (cf. §) (33-01b) `root, trunk', Tib. {\tib{ཁུལ་མ་}} \emph{khul-ma}
\textless{} *kurl (cf. §) `bottom or side of sth'
\item
Chi. 銀 \emph{ngin} \textless{} *ŋrə[n] \textless{} *ŋrərl (cf. §)
(33-01k) `silver,' Tib. {\tib{དངུལ་}} \emph{dṅul} \textless{} *dṅurl (cf. §)
'silver', OBur. {\burm{ငုယ်}} \emph{ṅuy} \textless{} *ṅurl (cf. §) `silver'
\item
Chi. 齦 \emph{ng{\sil{jɨ}}n} \textless{} *ŋə[n] \textless{} *ŋərl (cf. §)
(33-01-) `gums', Tib. 
%{\protect\[hypertarget{anchor-48}{}{}
Chi. 齦 \emph{ng{\sil{jɨ}}n} \textless{} *ŋə[n] \textless{} *ŋərl (cf. §) (33-01-)
'gums', Tib. {\tib{རྙིལ་}} \emph{rñil} \textless{} *rṅʸil \textless{} * rṅʸirl
(cf. §) `gums'
\item
Chi. 分 \emph{pjun} \textless{} *pə[n] \textless{} *pərl (cf. §)
(33-30a) `divide', Tib. {\tib{འབུལ་}} \emph{ḫbul }\textless{} *ḫburl, 
{\tib{འཕུལ་}} \emph{ḫphul} \textless{} *ḫphurl (cf. §) `give'
\item
Chi. 塵 \emph{drin} \textless{} *[d]rə[n] \textless{}
*[d]rərl (cf. §) (33-17a) `dust', Tib. {\tib{རྡུལ་}} \emph{rdul} \textless{}
*rdurl (cf.~§) `dust'
\item
Chi. 貧 \emph{bin} \textless{} *(Cə.){[}b]rə[n] \textless{}
*(Cə.){[}b]rərl (cf. §) (33-30v) `poor', Tib. {\tib{དབུལ་}} \emph{dbul}
\textless{} *dburl (cf. §)'
\item
Chi. 軍 \emph{kjun} \textless{}
*[k]{\sil{ʷ}}ər \textless{} *[k]{\sil{ʷ}}ərl (cf.§) (34-13a) `army', Tib.
 {\tib{གཡུལ་}} \emph{\textbf{g.yul}} \textless{} *gyurl
(cf. §) 'army, battle'
\end{description}

\paragraph{Inability to recognize -r- in
certain circumstances.} A wide variety of sound changes in the
history of the Chinese language lead to a merger of syllables with and
without medial *-r-. Thus, in many circumstance it is difficult to know
whether Old Chinese has a medial *-r-. Using `K' for velars, uvulars,
labio-velars, and labio-uvulars, `P' for labials, and `C' for any
consonants, the specific environments in which medial *-r- is lost
without trace are *K(r)ə, *K(r)u, *K(r)ək, *K(r)uk, *K(r)oŋ, *P(r)u,
*P(r)uk, *P(r)a, *Praj, *P(r)aw, and *C(r)ok in type B syllables, as
well as *C{\sil{ˤ}}(r)o in type A syllables (Baxter \& Sagart 2014: 214, 218,
243, 245).

§a. Some comparisons suggest a lack of medial *-r- in Old Chinese.

\begin{description} \item
Chi. 極 \emph{gik} \textless{} *[g](r)ək (05-04e) `ridge of roof,
highest point, centre', Tib. {\tib{མཁའ་}} \emph{mkhaḫ} `sky'

\item
Chi. 後 \emph{huwX} \textless{} *{\sil{[ɢ]ˤ}}(r)oʔ (10-08a) `behind', Tib. {\tib{འོག་}} \emph{ḫog} `under'

\item
Chi. 舅 \emph{gjuwX} \textless{} *[g](r)uʔ (04-16b) `maternal
uncle', Tib. {\tib{ཁུ་}} \emph{khu} `paternal uncle'

\item
Chi. 蜂�� \emph{phjowng} \textless{} *p{\sil{ʰ}}(r)oŋ (12-25st) `bee', Tib. {\tib{བུང་བ་}}
\emph{buṅ-ba}
\end{description}

§b. Other comparisons appear to confirm the presence of a medial *-r- in
Old Chinese.

\begin{description} \item
Chi. 蚼 \emph{xuwX} \textless{}{\sil{*qʰ}}{\sil{ˤ}}(r)oʔ (10-01n) `ant', Tib. {\tib{གྲོག་མོ་
}} \emph{grog-mo}

\item
Chi. 羽 \emph{hjuX} \textless{} *{\sil{ɢʷ}}(r)aʔ (01-24a) `feather', Tib. {\tib{སྒྲོ་
}} \emph{sgro} \textless{} *sg{\sil{ʷ}}ra

\item
Chi. 詖 \emph{pje} \textless{} *p(r)aj (18-16h) `one-sided, insincere
words', Tib. {\tib{ཕྲ་མོ་}} \emph{phra-mo} `slander'

\item
Chi. 披 \emph{phje }\textless{} *p{\sil{ʰ}}(r)aj (18-16j) `divide', Tib. {\tib{འབྲལ་
}} \emph{ḫbral} `be separate, to part', Bur. {\burm{ပြား}} \emph{prāḥ} `be divided
into parts'
\end{description}

\paragraph{Coblin's conjecture, *Tr- \textless{} *rT-}. Coblin
proposes that the correspondence of Chinese *Tr with Tibetan Tr- be
reconstructed *rT in the proto-language (1986: 21-22). Gong reiterates
this suggestion, further conjecturing that *rT continued into the Old
Chinese period (2002{[}2000]: 171). Handel concurs with Gong that *Tr
should take the form *rT in Old Chinese (2009: 211-217); he offers the
Chinese internal evidence that `dimidiated clusters' suggest *Kr- (e.g.
康良 \emph{khang ljang} `empty' \textless{} *{\sil{kʰ}}r-) and *Pr- (e.g. 朦朧
 \emph{muwng luwng} `indistinct' \textless{} *mr{\sil{ˤ}}-) but not *tr- (2009:
215). Because this proposal improves cor­re­spon­dences to Tibetan, I
assume it as a working hypothesis. The orthographic representation of
Coblin's conjecture as a change from the proto-language to Old Chinese
via a metathesis, i.e. *rT- \textgreater{} *Tr- does not necessarily
have its face value. Instead, *‑r- in Baxter \& Sagart's system indexes
all origins of Middle Chinese retroflex initials (cf. §), division-ii
rimes (cf. §), and \emph{chóngniǔ} rank-iii rimes (cf. §), including
both medial *-r- and, if one accepts Coblin's conjecture, initial *r-,
at least before dentals. In other words, if Coblin's conjecture is
correct then the sequence *rT- existed as such in Old Chinese.

\begin{description} \item
Chi. 撞 \emph{draewng} \textless{} *[N-t]{\sil{ˤ}}roŋ \textless{}
*[N-]r{[}t]{\sil{ˤ}}oŋ (12-08f') `strike', Tib. {\tib{རྡུང་}} \emph{rduṅ}
\item
Chi. 塵 \emph{drin} \textless{} *[d]rə[n] \textless{}
*r{[}d]ə[n] (33-17a) `dust (n.)', Tib. {\tib{རྡུལ་}} \emph{rdul}
\item
Chi. 冢 \emph{trjowngX} \textless{} *[t]roŋʔ \textless{}
*r{[}t]oŋʔ (11-13h), `tomb mound', Tib. {\tib{རྡུང་}} \emph{rduṅ} `small
mound, hillock', WBur. {\burm{တောင်}} \emph{toṅ} \textless{} *tuṅ `hill,
mountain'
\item
Chi. 晝 \emph{trjuwH} \textless{} *truks \textless{} *rtuks (14-09a)
'time of daylight', Tib. {\tib{གདུགས་}} \emph{gdugs} `mid-day, noon'
\item
Chi. 綴 \emph{trjwet} \textless{} *trot \textless{} *rtot (22-10b)
'bind', Tib. √rtod (pres. {\tib{རྟོད་}}) `tether, fasten, secure'
\item
Chi. 椓 \emph{traewk} \textless{} *tr{\sil{ˤ}}ok \textless{} *rt{\sil{ˤ}}ok (11-13c)
'beat, strike', Tib. {\tib{རྡུག་}} \emph{rdug} `conquer'
\end{description}

§a. Coblin's conjecture still leaves many examples of medial *-r- in
Chinese without a corresponding rhotic in Tibetan or Burmese.

\begin{description} \item
Chi. 殼 \emph{khaewk} \textless{} *[{\sil{kʰ}}]{\sil{ˤ}}rok (11-03a) `hollow shell,
hollow', Tib. {\tib{སྐོག་}} \emph{skog} `shell, peel', WBur. {\burm{ခေါက်}} \emph{khok}
\textless{} *ˀkuk `bark (n.)', Atsi  \emph{ˀku⁵⁵}
\item
Chi. 銀 \emph{ngin} \textless{} *ŋrə[n] (33-01k) `silver,' Tib.{\tib{དངུལ་}} \emph{dṅul}, OBur. {\burm{ငုယ်}} \emph{ṅuy}
\item
Chi. 眉 \emph{mij} \textless{} *mrər (27-14a) `eyebrow', WBur. {\burm{မွေး}} \emph{mweḥ} \textless{} *muyḥ `body hair'
\item
Chi. 貧 \emph{bin} \textless{} *(Cə.){[}b]rə[n] (33-30v) `poor',
Tib. {\tib{དབུལ་}} \emph{dbul}
\item
Chi. 虎 \emph{xuX} \textless{} {\sil{*qʰ}}{\sil{ˤ}}raʔ `tiger' (01-18b), Tib.
\protect\hypertarget{anchor-49}{}{}Chi. 虎 \emph{xuX} \textless{}
{\sil{*qʰ}}{\sil{ˤ}}raʔ `tiger' (01-18b), Tib. {\tib{སྟག་}} \emph{stag}
\item
Chi. 膚 \emph{pju} \textless{} *pra (01-51g) `skin',\footnote{Perhaps
  cognate to the {\burm{
  ပြား}} \emph{prāḥ} (\textless{} *brāḥ) of WBur. {\burm{အရေပြား}} \emph{a-re-prāḥ} `skin' in which case the Tibetan form is potentially
  not cognate.} Tib. {\tib{ལྤགས་}} \emph{lpags}
\item
Chi. 蝨 \emph{srit} \textless{} *sri{[}k] (29-35a) `louse', Tib. {\tib{ཤིག་}} \emph{śig}
\item
Chi. 殺 \emph{sreat} \textless{} *srat (21-29d) `kill', Tib. √sad (pres. {\tib{གསོད་}} \emph{gsod}), Bur. {\burm{သတ်}} \emph{sat} \textless{} *ˀsat, Lashi
 \emph{ˀsa:tH}
 \item
Chi. 三 \emph{sam} \textless{} *srum (38-30a) `three', Tib. {\tib{གསུམ་}} \emph{gsum}, Bur. {\burm{သုံး}} \emph{suṃḥ} \textless{} *suṃḥ, Lashi  \emph{sɔmH}
\item
Chi. 沙 \emph{srae} \textless{} *s{\sil{ˤ}}raj (18-15a) `sand', Tib. {\tib{ས་}} \emph{sa} `earth', Bur. {\burm{သဲ}} \emph{sai} `sand'
\end{description}

§b. Handel offers an explanation for the correspondence of *sr- in
Chinese to Tibetan \emph{s}- (or \emph{ś}-) seen in the immediately
preceding four examples, namely that Chinese maintains inherited *sr-,
which Tibetan (and all languages of the family except Chinese) simplify
to \emph{s}- before vowels other than inherited *-e- and *-i- (2009:
201-209). He discusses the following examples.\footnote{Handel
  additionally compares Chi. 鼪~ \emph{sjengH} \textless{} *sreŋ (09-25u)
  `weasel' to Tibetan {\tib{སྲེ་མོང་}} \emph{sre-moṅ} `mongoose' (2009: 201), a
  comparison he derives from Benedict (1972: 79). These forms are not a
  good phonetic match. In addition, mongooses are not native to Tibet,
  so if the words are connected it is not as cognates. Handel
  additionally mentions two comparison noted by Coblin (1986: 115, 128),
  viz. Chi. 雙  \emph{sraewng} \textless{} *s{\sil{ˤ}}roŋ `a pair' (12-24a), Tib. {\tib{ཟུང་}} \emph{zuṅ} \textless{} *dzuŋ, Bur.{\burm{စုံ}} -\emph{cuṃ} and Chi. 率
   \emph{srwijH} \textless{} *s-rut-s (31-23a) `lead (v.); commander',
  Chi. 帥 \emph{srwijH} \textless{} *s-rut-s (31‑24a) `leader (of an
  army)', Tib. {\tib{སྲིད་}} \emph{srid} `government'. The correspondence seen
  in `pair' conforms to Handel's proposal in so far as the Tibetan word
  lacks a medial -r-, but it contradicts his proposal in that the
  Tibetan initial is z- \textless{} *dz- rather than s-. Jacques rejects
  this pair as cognates because of the irregular initial correspondence
  (2015: 218 note 18 on p. 221). Because the second comparison shows an
  irregular vowel correspondence, it is difficult to evaluate whether it
  supports or contradicts Handel's proposal. One may note that no
  Tibetan words begin *sru-, and consequently suggest that Tib. {\tib{སྲིད་
 }} \emph{srid} `government' derives regularly from *srud. Jacques (2015:
  216 note 2 on p. 221) rejects this comparison, both on the basis of
  the irregular vowel correspondences, and also because he sees the
  Chinese word as deriving from 律 \emph{lwit} \textless{} *rut
  (31-18c). Note that Tib. {\tib{སྲིད་}} \emph{srid} does not mean `government'
  in isolation but only in the compounds 
  {\tib{ཆབ་སྲིད་}} \emph{chab-srid} or
  {\tib{ཆུ་སྲིད་}} \emph{chu-srid}
  composed of `water' plus `possible, exist'
  (cf. Ishikawa 1998).}

\begin{description} \item
Chi. 殺  \emph{sreat \textless{} }*srat (21-29d) `kill', Tib. √sad (pres. {\tib{གསོད་}} \emph{gsod}), Bur. {\burm{သတ်}} \emph{sat} \textless{} *ˀsat, Lashi
 \emph{ˀsa:tH}
\item
Chi. 沙  \emph{srae} \textless{} *s{\sil{ˤ}}raj (18-15a) `sand', Tib. {\tib{ས་}} \emph{sa} `earth', Bur. {\burm{သဲ}} \emph{sai} `sand'
\item
Chi. 蝨  \emph{srit} \textless{} *sri{[}k] (29-35a) `louse', Tib. {\tib{ཤིག་
}} \emph{śig}
\item
Chi. 甥  \emph{sraeng} \textless{} *s.reŋ (09-25g) `son-in-law', Tib.
{\tib{སྲིང་མོ་}} \emph{sriṅ-mo} `sister of a man'
\item Chi. 產  \emph{sreanX} \textless{} *s-ŋrarʔ (34-46a) `produce', Tib.
{\tib{སྲེལ་}} \emph{srel} `rear, bring up'
\item
Chi. 欶  \emph{sraewk} \textless{} *s{\sil{ˤ}}rok (11-21o) `suck, inhale, WBur. {\burm{သောက်}} \emph{sok} \textless{} *śuk `drink', Lashi  \emph{śu:kH}
\end{description}

To these one can further offer the following comparisons:

\begin{description} \item
Chi. 三  \emph{sam} \textless{} *srum (38-30a) `three', Tib. {\tib{གསུམ་
}} \emph{gsum, }Bur. {\burm{သုံး}} \emph{suṃḥ} \textless{} *suṃḥ, Lashi  \emph{sɔmH}

\item
Chi. 獺  \emph{that} \textless{} *r̥{\sil{ˤ}}at (21-24i) `otter', Tib.
\protect\hypertarget{anchor-50}{}{}Chi. 獺  \emph{that} \textless{} *r̥{\sil{ˤ}}at
(21-24i) `otter', Tib. {\tib{སྲམ་}} \emph{sram}, Bur. {\burm{ဖျံ}} \emph{phyaṃ}, Lashi
 \emph{ˀśam}
\end{description}

The five examples `kill', `sand', `three', `son-in-law', and `produce'
conform to the pattern Handel describes. He acknowledges that `louse'
con­tra­dicts his hypothesis, finding it best to ``attribute the
variation in Tibetan to an unknown cause'' (2009: 209) and criticizing
Benedict (1972: 108 n. 304) for setting up *śr- and *sr- on the sole
basis of the two outcomes *śr- \textgreater{} \emph{ś}- and *sr-
\textgreater{} \emph{sr}- in Tibetan. The same criticism applies to
Jacques' reconstructions of Tibetan \emph{ś}- \textless{} *sr- versus
 \emph{sr}- \textless{} *sə.r- to explain the same two examples (2015:
217). The different re­con­structed initials *sr- versus *s.r- in Baxter
\& Sagart's system implies an explanation similar to Jacques'. For
'otter' Handel resolves the obstacle of \emph{sr}- in Tibetan appearing
before -\emph{a}- by reconstructing *s-ram `otter' rather than *sram;
for him there is ``no apparent cognate in Chinese'' (2009: 202 note 20).
This *sr- versus *s-r- notational slight of hand, exactly akin to the
tack of Benedict that Handel rejects, has no descriptive or explanatory
power.\footnote{From these scant data in support of a Tibetan innovation
  (i.e. *sr- \textgreater{} s- except before -e- and i-), Handel
  concludes that ``the loss of medial *-r- after initial *s- in PTB
  {[}Proto-Tibeto-Burman] is a common innovation that provides
  evidence supporting the early split of Chinese from the common PST
  stock, before the subsequent ramification of Tibeto-Burman'' (Handel
  2009: 204). This reasoning would be correct if Benedict's
  reconstructions reflected systematic correspondences among all the
  so-called `Tibeto-Burman' languages, but they do not.}

§c. Guillaume Jacques rejects Handel's explanation of the development of
*sr-. First, he objects to the comparison of Chinese 產  \emph{sreanX}
\textless{} *s-ŋrarʔ (34-46a) `produce' to Tibetan 
{\tib{སྲེལ་}} \emph{srel}
'rear, bring up' (2015: 216 note 2 on p. 221), but in so doing puts
undue weight on Baxter \& Sagart's reconstruction. Handel reconstructs
*s{\sil{ˤ}}renʔ (2009: 200, 204) and Schuessler *s{\sil{ˤ}}ranʔ / *s{\sil{ˤ}}renʔ (2009: 291).
Second, Jacques points out that Japhug Rgyalrong maintains medial -r- in
two cognates where Chinese has the vowel *ə.

\begin{description} \item
Chi. 色  \emph{srik} \textless{} *srək (05-31a) `color, sex, shame',
Japhug Rgy. \emph{tɤ‑zraʁ} `shame'

\item
Chi. 參 \emph{syim}\textless{} *srəm (38-29a) `rhizome', Japhug Rgy.
 \emph{tɤ-zrɤm} `root'\footnote{Jacques (2017c) further proposes Tib. {\tib{གཤམ་}} \emph{gśam} `base, underpart' as a cognate. }
\end{description}

These two examples ensure that Handel's proposal that all languages
except Chinese change *sr- to \emph{s}- before all vowels except *-e-
and *-i- cannot be correct. Jacques (2015: 219-220) offers three
explanations for the attested correspondences: (1). Languages other than
Chinese share a change *sr- to \emph{s}-, conditioned be a following
inherited *-a- and possibly also *-o-. (2). Chinese innovates medial
*-r- through infixation (cf. Baxter \& Sagart 2014: 57-58). (3). Current
re­con­structions of Old Chinese include too many instances of medial
*-r-. Jacques remains agnostic among these three possibilities.

\paragraph{Pulleyblank's conjecture, *Cr- \textless{} *RC-}.
Returning to the Tibetan correspondences of those Chinese words that
have an unexpected medial *-r- permits further contemplation of Jacques'
third option. In most cases Tibetan exhibits a complex onset with a {\tib{སྔོན་འཇུག་}} \emph{sṅon-ḥjug} or {\tib{མགོ་ཅན་}} \emph{mgo-can} consonant, i.e.
 \emph{s}-,  \emph{d},  \emph{l}-, or \emph{g}-. The overall pattern of
unexpected Chinese *-r- cor­re­sponding to Tibetan complex onsets
entitles one to speculate that Chinese had some `pre-initial' in these
words that conditioned the same effects in Middle Chinese as Old Chinese
*-r-. This speculation extends Coblin's conjecture to environments
beyond the dental stops. Because it is unlikely that specifically *r-
was the `pre-initial' in all of these words, it is more suitable to
reconstruct *R- as a convention to mean `an indeterminate initial of a
complex onset which bears the same effects as *-r-'. Particularly based
on the comparison of Tibetan {\tib{གསོད་}} \emph{gsod} `kill' with Chinese 殺
 \emph{sreat} \textless{} *srat (21-29d), Pulleyblank posits *ks- as an
origin of Middle Chinese 生 \emph{sr}- (1965: 206-207). Gong similarly
proposes *rsat as the origin of Chinese 殺 \emph{sreat} \textless{}
*srat (21-29d) (2002{[}2001]: 171).

\begin{description} \item
Chi. 殼  \emph{khaewk} \textless{} *[{\sil{kʰ}}]{\sil{ˤ}}rok \textless{}
*R{[}{\sil{kʰ}}]{\sil{ˤ}}ok (11-03a) `hollow shell, hollow', Tib. {\tib{སྐོག་}} \emph{skog}
'shell, peel'
\item
Chi. 銀  \emph{ngin} \textless{} *ŋrə[n] \textless{} *Rŋə[n]
(33-01k) `silver', Tib. {\tib{དངུལ་}} \emph{dṅul}
\item
Chi. 貧  \emph{bin} \textless{} *(Cə.){[}b]rə[n] \textless{}
*(Cə.)R{[}b]ə[n] (33-30v) `poor', Tib. {\tib{དབུལ་}} \emph{dbul}
\item
Chi. 虎  \emph{xuX} \textless{} {\sil{*qʰ}}{\sil{ˤ}}raʔ `tiger' \textless{} *R{\sil{qʰ}}{\sil{ˤ}}aʔ
(01-18b), 
Tib. 
%{\protect\[hypertarget{anchor-51}{}{}
Chi. 虎 \emph{xuX}
\textless{} {\sil{*qʰ}}{\sil{ˤ}}raʔ `tiger' \textless{} *R{\sil{qʰ}}{\sil{ˤ}}aʔ (01-18b), Tib. {\tib{སྟག་}} \emph{stag}
\item
Chi. 膚  \emph{pju} \textless{} *pra \textless{} *Rpa (01-51g) `skin',
Tib. {\tib{ལྤགས་}} \emph{lpags} (p. n. )
\item
Chi. 殺  \emph{sreat} \textless{} *srat \textless{} *Rsat (21-29d)
'kill', Tib. √sad (pres. {\tib{གསོད་}} \emph{gsod})
\item
Chi. 三  \emph{sam} \textless{} *srum \textless{} *Rsum (38-30a) `three',
Tib. {\tib{གསུམ་}} \emph{gsum}
\end{description}

§a. There are two comparisons, in which the Tibetan cognates have a
simplex onset; these are apparent counterexamples to Pulleyblank's
conjecture.

\begin{description} \item
Chi. 蝨  \emph{srit} \textless{} *sri{[}k] (29-35a) `louse', Tib. {\tib{ཤིག་}} \emph{śig}
\item Chi. 沙  \emph{srae} \textless{} *s{\sil{ˤ}}raj (18-15a) `sand', Tib. {\tib{ས་}} \emph{sa} `earth'
\end{description}

In the first of these apparent exceptions, the Japhug Rgyalrong cognate
 \emph{zrɯɣ} `louse' confirms the Chinese onset *sr- is original. One may
suggest that Tibetan underwent a change *sr- \textgreater{} ś-
conditioned by the vowel -i-. Because the change *sr- \textgreater{} ś-
preceded Dempsey's law *-eŋ \textgreater{} -\emph{iṅ} (cf. §), Tibetan {\tib{སྲིང་མོ་}} \emph{sriṅ-mo} \textless{} *sreṅ-mo `sister of a man' avoided
the former change (cf. Chi. 甥 \emph{sraeng} \textless{} *s.reŋ
{[}09-25g] `son-in-law'). I have no explanation for the correspondence
seen in `sand'.

§b. Fitting Burmese evidence into the overall picture of the sibilants
is not straightforward (cf. §).

\paragraph{Old Chinese *-k \textless{} *-k, *-kə}.
Karlgren (1931: 43-44) noticed that Old Chinese *-k often corresponds to
an open syllable in Tibetan. Karlgren consequently suggests that Chinese
*-k is a `suffix', although endowed with no morphological significance.
Although Burmese always has an open syllable in the relevant cases (cf.
§), in many Tibetan has a final -\emph{ḫ} {[}-x] (cf. Hill 2011
'inventory': 453 and cf. §). The correspondence of Chinese *-k to velar
stop finals in Tibetan and Burmese must be reconstructed *-k in the
common ancestor (cf. §a). The correspondence to open syllables (and
Tibetan -\emph{ḫ)} is more difficult to reconstruct. It is possible to
project the Tibetan {[}-x] back on to the proto-language, but this
proposal suggests an unlikely fortition in the history of Chinese *-x
\textgreater{} *‑k. Thai underwent a similar fortition in syllable
onsets (Brown 1956: 62, 146; Li 1977: 207-214), but I am unaware of a
case of unconditioned fricative hardening in syllable final position. A
phonetically more plausible reconstruction is *-kə, developing to *-k in
Chinese via apocope, and undergoing a series of changes such as *-kə
\textgreater{} *-gə- \textgreater{} *-ɣə \textgreater{} *-ɣ
\textgreater{} *‑x in Tibetan and Burmese, with subsequent unconditioned
loss of *-x in Burmese and the same loss in Tibetan during historical
times, conditioned by all vowels except for -\emph{a}.

§a. Examples to be reconstructed *-k :
\begin{description} \item
Chi. 夜  \emph{yaeH} \textless{} *N.raks (02-27j) `night', Tib. {\tib{ཞག་}} \emph{źag} \textless{} *rʸak (Benedict's law, cf. §) `day, 24hrs', Bur. {\burm{ရျက်}} \emph{ryak}
\item
Chi. 隻  \emph{tsyek} \textless{} *tek (08-11c) `single', Tib. {\tib{གཅིག་}} \emph{gčig} \textless{} *gtʸig `one', Bur. {\burm{တစ်}} \emph{tac} \textless{}
*dik (Wolfenden's law, cf.~§)
\item
Chi. 燭  \emph{tsyowk }\textless{} *tok (11-12e) `torch', Tib. {\tib{དུགས}}
'kindle, light (v.), WBur. {\burm{တောက်}} \emph{tok} \textless{} *duk (Maung
Wun's law, cf. §) `blaze, flame, shine'
\item
Chi. 毒  \emph{dowk} \textless{} *d{\sil{ˤ}}uk (14-05a) `poison', Tib.
\protect\hypertarget{anchor-52}{}{}Chi. 毒 \emph{dowk} \textless{} *d{\sil{ˤ}}uk
(14-05a) `poison', Tib. {\tib{དུག་}} \emph{dug}, WBur. {\burm{တောက်}} \emph{tok}
\textless{} *duk (Maung Wun's law) `be toxic'
\item
Chi. 六  \emph{ljuwk} \textless{} *k.ruk (14-16a) `six', Tib. {\tib{དྲུག་
}} \emph{drug}, WBur. {\burm{ခြောက်}} \emph{khrok} \textless{} *kruk (Maung Wun's
law), Lashi  \emph{khjukH}
\end{description}

For a complete list of examples of this correspondence see §.

§b. Examples to be reconstructed *-kə :

\begin{description} 
\item
Chi. 渡  \emph{duH} \textless{} *d{\sil{ˤ}}ak-s \textless{} *d{\sil{ˤ}}akə-s\footnote{One
  may assume that the suffix -s was only applied after the loss of *-ə,
  but the current notation is most explicit and therefore most helpful.}
(02-16b) `ford', Tib. {\tib{འདའ་}} \emph{ḫdaḫ} `pass (v.)'

\item
Chi. 射  \emph{zyek} \textless{} *Cə.lak (02-26a) `hit with bow and
arrow', Tib. {\tib{མདའ་}} \emph{mdaḫ} \textless{} *mlaḫ (Bodman's law, cf. §)
\textless{} *mlakə `arrow', OBur. {\burm{မ္လား}} \emph{mlāḥ}
\item
Chi. 澤  \emph{draek} \textless{} *l{\sil{ˤ}}rak \textless{} *l{\sil{ˤ}}rakə (02-25o)
`marsh, moisture', Tib. %{\protect\[hypertarget{anchor-53}{}{}
Chi. 澤
 \emph{draek} \textless{} *l{\sil{ˤ}}rak \textless{} *l{\sil{ˤ}}rakə (02-25o) `marsh,
moisture', Tib. {\tib{བཞའ་}} \emph{bźaḫ }\textless{} *blʸaḫ (Benedict's law, cf.
§) `wet'
\item
Chi. 百 \emph{paek} \textless{} *p{\sil{ˤ}}rak \textless{} *p{\sil{ˤ}}rakə (02-37a)
'hundred', Tib. {\tib{བརྒྱ }} \textless{} OTib. {\tib{བརྒྱའ་}} \emph{brgyaḫ} (PT 1111, l.
5 \emph{et passim}) \textless{} *bryaḫ (Li's law, cf. §) \textless{}
*bryakə, OBur. {\burm{ရျာ}} \emph{ryā}
\item
Chi. 魄  \emph{phaek} \textless{} *p{\sil{ʰ}}{\sil{ˤ}}rak \textless{} *p{\sil{ʰ}}{\sil{ˤ}}rakə (02-38o)
'soul', Tib. {\tib{བླ་}} \emph{bla} \textless{} OTib. {\tib{བརླའ་}} \emph{brlaḫ} (PT
1287, l. 57), WBur. {\burm{လိပ်ပြာ}} \emph{lipprā} `butterfly, soul'
\item
Chi. 極  \emph{gik} \textless{} *[g](r)ək \textless{} *[g](r)əkə
(05-04e) `ridge of roof, highest point, centre', Tib. {\tib{མཁའ་}} \emph{mkhaḫ}
\textless{} *mkakə `sky'
\item
Chi. 易  \emph{yek} \textless{} *lek \textless{} *lekə `change; exchange'
(08-12a), Tib. {\tib{རྗེ་}} \emph{rǰe} \textless{} *rlʸe `exchange', Bur. {\burm{လဲ}} \emph{lai} `exchange' (cf. p. n. )
\item
Chi. 角  \emph{kaewk} \textless{} *C.{\sil{kˤ}}rok \textless{} *C.{\sil{kˤ}}rokə (11-02a)
'horn, corner', Tib. {\tib{རུ་}} \emph{ru} `horn', {\tib{གྲུ་}} \emph{gru} `corner',
WBur. {\burm{ချို}} \emph{khyui} `horn'

\item Chi. 蝮  \emph{phjuwk} \textless{} *p{\sil{ʰ}}uk \textless{} *p{\sil{ʰ}}ukə (14-26j) `a
kind of snake', Tib. {\tib{འབུ་}} \emph{ḫbu} `worm, insect', OBur. {\burm{ပိုဝ်း
}} \emph{puiwḥ} `insect'

Chi. 日  \emph{nyit} \textless{} *C.ni{[}k] \textless{} *C.nikə
(29-26a) `sun', Tib. {\tib{ཉི་མ་}} \emph{ñi-ma} `sun' \textless{} *nʸi, OBur. {\burm{နိယ်}} \emph{niy} `sun',

\item
Chi. 桼 \emph{tshit} \textless{} *ts{\sil{ʰ}}i{[}k] \textless{} *ts{\sil{ʰ}}ikə
(29-32a) `varnish', Tib. %{\protect\[hypertarget{anchor-54}{}{}
Chi. 桼
 \emph{tshit} \textless{} *ts{\sil{ʰ}}i{[}k] \textless{} *ts{\sil{ʰ}}ikə (29-32a)
'varnish', Tib. {\tib{ཚི་}} \emph{tshi} `sticky matter', WBur. {\burm{စေး}} \emph{ceḥ}
\textless{} OBur. {\burm{*ciyḥ}} `be sticky, adhesive'
\end{description}

\paragraph{Old Chinese *-ʔ \textless{} *ʔ, *-q}. Chinese final
glottal stop sometimes corresponds to a -\emph{k} in Tibetan or Burmese
(Bodman 1980: 134 et passim, Zhèngzhāng 2000: 66, cf. §§, ). Tibetan and
Burmese words ending in velar stops most frequently correspond to
Chinese *-k (cf. §) and Chinese *-ʔ also correspond to Tibetan and
Burmese open syllables. To account for this triple correspondence (Chi.
*-k :: Tib. -\emph{g} :: Bur. -\emph{k}; Chi. *-ʔ :: Tib. -\emph{g} ::
Bur. -\emph{k}; Chi *-ʔ :: Tib. Ø :: Bur. -Ø) it is necessary to
reconstruct a Trans-Himalayan segment *-q that becomes *-ʔ in Chinese
but -k in Tibetan and Burmese. Sagart (2017) points out that is favor of
this reconstruction are such Austronesian cognates as *punuq `brain'.

§a. Examples to be reconstructed *-k :

\begin{description} \item
Chi. 夜  \emph{yaeH} \textless{} *N.raks (02-27j) `night', Tib. {\tib{ཞག་
}} \emph{źag} \textless{} *rʸak `day, 24hrs', Bur. {\burm{ရျက်}} \emph{ryak}
\item
Chi. 隻  \emph{tsyek} \textless{} *tek (08-11c) `single', Tib. {\tib{གཅིག་}} \emph{gčig} \textless{} *gtʸig, Bur. {\burm{တစ်}} \emph{tac} \textless{} *dik
'one'
\item
Chi. 燭  \emph{tsyowk }\textless{} *tok (11-12e) `torch', Tib. {\tib{དུགས}}
'kindle, light (v.), WBur. {\burm{တောက်}} \emph{tok} \textless{} *duk `blaze,
flame, shine'
\item
Chi. 毒  \emph{dowk} \textless{} *d{\sil{ˤ}}uk (14-05a) `poison', Tib.
\protect\hypertarget{anchor-55}{}{}Chi. 毒 \emph{dowk} \textless{} *d{\sil{ˤ}}uk
(14-05a) `poison', Tib. {\tib{དུག་}} \emph{dug}, WBur. {\burm{တောက်}} \emph{tok}
\textless{} *duk `be toxic'
\item
Chi. 六  \emph{ljuwk} \textless{} *k.ruk (14-16a) `six', Tib. {\tib{དྲུག་}} \emph{drug}, WBur. {\burm{ခြောက်}} \emph{khrok} \textless{} *kruk `six', Lashi
 \emph{khjukH}
\end{description}

For a \emph{ }complete list of examples of this correspondence see §.

§b. Examples to be reconstructed *-q :

\begin{description} \item
Chi. 語  \emph{ngjoX} \textless{} *ŋaʔ \textless{} *ŋaq (01-29t) `speak',
Tib. {\tib{ངག་}} \emph{ṅag} `speech'

\item
Chi. 女  \emph{nrjoX} \textless{} *nraʔ \textless{} *nraq (01-56a)
'woman', Tib. {\tib{ཉག་མོ་}} \emph{ñag-mo} (cf. p. n.~)

\item
Chi. 武  \emph{mjuX} \textless{} *maʔ \textless{} *maq (01-71a)
'military', Tib. {\tib{དམག་}} \emph{dmag} `army'

\item
Chi. 友  \emph{hjuwX} \textless{} *{\sil{ɢʷ}}əʔ \textless{} *{\sil{ɢʷ}}əq (04-17e)
'friend', Tib. {\tib{གྲོགས་}} \emph{grogs} \textless{} *g{\sil{ʷ}}raks

\item
Chi. 婦  \emph{bjuwX} \textless{} *mə.bəʔ \textless{} *mə.bəq (04-63a)
`woman, wife', Tib. 
%{\protect\[hypertarget{anchor-56}{}{}
Chi. 婦
 \emph{bjuwX} \textless{} *mə.bəʔ \textless{} *mə.bəq (04-63a) `woman,
wife', Tib. {\tib{བག་མ་}} \emph{bag-ma} `bride'
\item
Chi. 後  \emph{huwX} \textless{} *{\sil{[ɢ]ˤ}}(r)oʔ \textless{}
*{\sil{[ɢ]ˤ}}(r)oq (10-08a) `behind', Tib. {\tib{འོག་}} \emph{ḫog} `under'
\item
Chi. 嘔 \emph{'uwX} \textless{} {\sil{*qˤ}}(r)oʔ \textless{} {\sil{*qˤ}}(r)oq (10-10i)
'vomit', 
%Tib. {\protect\hypertarget{anchor-57}{}
%Chi. 嘔 `\emph{uwX}
%\textless{} {\sil{*qˤ}}(r)oʔ \textless{} {\sil{*qˤ}}(r)oq (10-10i) `vomit', 
Tib. 
{\tib{སྐྱུག་}} \emph{skyug}
\item
Chi. 蚼  \emph{xuwX} \textless{}{\sil{*qʰ}}{\sil{ˤ}}(r)oʔ \textless{}{\sil{*qʰ}}{\sil{ˤ}}(r)oq (10-01n)
'ant', Tib. {\tib{གྲོག་མོ་}} \emph{grog-mo}
\item
Chi. 帚  \emph{tsyuwX} \textless{} *[t.p]əʔ \textless{} *[t.p]əq
(13-20Aa) `broom', Tib. %{\protect\[hypertarget{anchor-58}{}{}
Chi. 帚
 \emph{tsyuwX} \textless{} *[t.p]əʔ \textless{} *[t.p]əq
(13-20Aa) `broom', Tib. {\tib{ཕྱགས་མ་}} \emph{phyags-ma}
\item
Chi. 腦  \emph{nawX} \textless{} *n{\sil{ˤ}}uʔ \textless{} *n{\sil{ˤ}}uq (16-28f)
'brain', WBur. {\burm{နှောက်}} \emph{nhok} \textless{} *ˀnuk
\end{description}

§c. Examples to be reconstructed *-ʔ :

\begin{description} \item
Chi. 苦  \emph{khuX} \textless{} *{\sil{kʰ}}{\sil{ˤ}}aʔ (01-01u) `bitter', Tib. {\tib{ཁ་
}} \emph{kha}, Bur. {\burm{ခါး}} \emph{khāḥ} \textless{} *kāḥ, Lashi  \emph{kho:H}

\item
Chi. 羽  \emph{hjuX} \textless{} *{\sil{ɢʷ}}(r)aʔ (01-24a) `feather', Tib. {\tib{སྒྲོ་
}} \emph{sgro} \textless{} *sg{\sil{ʷ}}ra

\item
Chi. 五 \emph{nguX} \textless{} *C.ŋ{\sil{ˤ}}aʔ (01-29a) `five', Tib. {\tib{ལྔ་
}} \emph{lṅa}, Bur. {\burm{ငါး}} \emph{ṅāḥ}

\item
Chi. 睹  \emph{tuX} \textless{} *t{\sil{ˤ}}aʔ (01-38c') `see', Tib. {\tib{ལྟ་
}} \emph{lta} `look at'

\item
Chi. 旅 \emph{ljoX} \textless{} *raʔ (01-55a) `military unit', Tib. {\tib{དགྲ་
}} \emph{dgra} `enemy'

\item
Chi. 咀  \emph{dzjoX} \textless{} *dzaʔ (01-57u) `eat', Tib. {\tib{ཟ་}} \emph{za}
\textless{} *dza, Bur. {\burm{စား}} \emph{cāḥ} \textless{} *dzāḥ, Lashi
 \emph{tsɔ:}

\item
Chi. 父  \emph{bjuX} \textless{} *N-paʔ (01-67a) `father', Tib. {\tib{ཕ་
}} \emph{pha}
\end{description}

§d. A few potential examples show problems of one kind or another.

\begin{description} \item
Chi. 虎  \emph{xuX} \textless{} {\sil{*qʰ}}{\sil{ˤ}}raʔ `tiger' (01-18b), Tib.
\protect\hypertarget{anchor-59}{}{}Chi. 虎 \emph{xuX} \textless{}
{\sil{*qʰ}}{\sil{ˤ}}raʔ `tiger' (01-18b), Tib. {\tib{སྟག་}} \emph{stag}

\item
Chi. 膚  \emph{pju} \textless{} *pra (01-51g) `skin', Tib. {\tib{ལྤགས་
}} \emph{lpags }(p. n. )

\item
Chi. 汝  \emph{nyoX} \textless{} *naʔ (01-56j) `you', Bur. {\burm{နင်}} \emph{naṅ}
\end{description}

Beckwith \& Kiyose reconstruct the Old Chinese word for `tiger' as
*staɣ, corresponding to *s.taʔ with Baxter \& Sagart's conventions. In
support of this reconstruction, Beckwith \& Kiyose point to the Japanese
word \emph{tora} `tiger', written 刀良 in the 正倉院戸籍帳 \emph{Shōsōin
Kosekichō} (702 C.E.). They rea­son­ably point out that Tibetan 
{\tib{སྟག་}} \emph{stag} and Japanese \emph{tora} look quite easy to relate and if
there is a relationship between them it must be by way of Chinese. In
addition, Beckwith points to a  \emph{xiéshēng} connection with 處
 \emph{tsyhoX} \textless{} *t.{\sil{qʰ}}aʔ (01-18a) `be at', noting that S.
Starostin re­con­structs this word *t{\sil{ʰ}}aʔ (2008: 9 note 50). I have no
explanation for the irregular cor­re­spondences of the second person
pronoun and the word `skin', except that the Tibetan word is potentially
not cognate (see p. n. ).

\paragraph{*-j \textless{} *-j, *-l}. Chinese merges what must
have been distinct *-j (cf. §a) and *-l (cf. §b) in the Trans-Himalayan
proto-language (see Hill 2014 `cognates').

§a. The correspondence of Chinese *-j with Burmese -\emph{y} points to a
shared inheritance of *‑j (cf. Schuessler 1974: 85-86), which Tibetan
lost (cf. §).

\begin{description} \item
Chi. 攲  \emph{kje} \textless{} *kraj (18-01d') `slanting', Chi. 掎
 \emph{kjeX} \textless{} *krajʔ (18-01y) `pull aside', Bur. {\burm{ကယ်
}} \emph{kay} `be distended'

\item
Chi. 嘉  \emph{kae} \textless{} *{\sil{kˤ}}raj (18-04g) `excellent', Bur. {\burm{ကဲ
}} \emph{kai} `overdo, exceed'

\item
Chi. 俄  \emph{nga} \textless{} *ŋ{\sil{ˤ}}aj (18-05h) `slanting', Bur. {\burm{ငဲ့
}} \emph{ṅaiʔ} `be inclined on one side'

\item
Chi. 㢋 \emph{tsyheX} \textless{} *k-l̥ajʔ (18-08t) `wide, extend', OBur. {\burm{က္လယ်}} \emph{klay} `wide, broad'

\item
Chi. 羅 \emph{la} \textless{} *r{\sil{ˤ}}aj (18-10a) `a kind of net', Tib. {\tib{དྲ་
}} \emph{dra} `net', Tib. {\tib{རྒྱ་}} \emph{rgya} \textless{} *rya (Li's law, cf.
§) `net, trap'

\item
Chi. 籬  \emph{lje} \textless{} *raj (18-11g) `hedge', Tib. {\tib{ར་}} \emph{ra}
'courtyard'

\item
Chi. 沙 \emph{srae} \textless{} *s{\sil{ˤ}}raj (18-15a) `sand', Tib. {\tib{ས་
}} \emph{sa} `earth', Bur. {\burm{သဲ}} \emph{sai} `sand'

\item
Chi. 波  \emph{pa} \textless{} *p{\sil{ˤ}}aj (18-16l) `wave', Tib. {\tib{དབའ་
}} \emph{dbaḫ} `wave' (cf. p. n. )

\item
Chi. 詖  \emph{pje} \textless{} *p(r)aj (18-16h) `one-sided, insincere
words', Tib. {\tib{ཕྲ་མོ་}} \emph{phra-mo} `slander'

\item
Chi. 跛 \emph{paX} \textless{} *p{\sil{ˤ}}ajʔ (18-16m) `walk lame', Bur. {\burm{ဖဲ
}} \emph{phai} `avoid, shun', Bur. {\burm{ဖယ်}} \emph{phay} `push aside'
\end{description}

A complete list of examples of this correspondence is given at §.

§b. The correspondence of Chinese *-j with Tibetan -\emph{l} points to a
change *l \textgreater{} *-j in the history of Chinese; Burmese
generally loses *l (cf. §), but keeps it as -\emph{y} after the vowel
-\emph{u}- (cf. §).

\begin{description} \item
Chi. 可  \emph{khaX} \textless{} *[k]{\sil{ʰ}}{\sil{ˤ}}a{[}j]ʔ (18-01a) `be able',
Chi. 可  \emph{gaX} \textless{} *gˁajʔ (18-01a) `carry', Tib. √kal (pres. {\tib{འགེལ་}} \emph{ḫgel}) `load', 
{\tib{འཁེལ་}} \emph{ḫkhel} `be loaded'

\item
Chi. 河  \emph{ha }\textless{} *C.{[}g]{\sil{ˤ}}aj (18-01g) `river', Tib. {\tib{རྒལ་
}} \emph{rgal} `cross, ford' (cf. p. n. )

\item
Chi. 荷 \emph{ha }\textless{} *[g]{\sil{ˤ}}aj (18-01o) `carry', Chi. 駕
 \emph{kaeH} \textless{} *kraj-s (18-04e) `yoke (v.)', Tib. {\tib{ཁལ་
}} \emph{khal} `burden, load (n.)', Bur. {\burm{က}} \emph{ka} `saddle'

\item
Chi. 歌 \emph{ka} \textless{} *{\sil{kˤ}}aj (18-01q) `sing, song', Bur. {\burm{က
}} \emph{ka} `dance (v.)'

\item
Chi. 加  \emph{kae }\textless{} *{\sil{kˤ}}raj (18-04a) `add', Tib. {\tib{བཀྲལ་
}} \emph{bkral} `appoint', Tib. {\tib{ཁྲལ་}} \emph{khral} `tax', Bur. {\burm{ကြား
}} \emph{krāḥ} `interval'

\item
Chi. 枷 \emph{kaeH} \textless{} *kraj-s (18-04c) `stand, support', Bur. {\burm{ကြာ}} \emph{krā} `last, take time'

\item
Chi. 義  \emph{ngjeH} \textless{} *ŋaj-s (18-05r) `duty, justice', Chi.
儀  \emph{ngje} \textless{} *ŋaj (18-05u) `proper demeanor, model', Bur. {\burm{ငှ}} \emph{ṅha} `distribute equally'

\item
Chi. 施 \emph{sye} \textless{} *l̥aj (18-09l') `give, bestow, extend',
Bur. {\burm{လျား}} \emph{lyāḥ} `extended, be prolonged'

\item
Chi. 離  \emph{lje} \textless{} *[r]aj (18-11f) `leave, distribute',
Tib. {\tib{རལ་}} \emph{ral} `rent, rift'

\item
Chi. 披  \emph{phje }\textless{} *p{\sil{ʰ}}(r)aj (18-16j) `divide', Tib. {\tib{འབྲལ་
}} \emph{ḫbral} `be separate, to part', Bur. {\burm{ပြား}} \emph{prāḥ} `be divided
into parts'

\item
Chi. 罷疲  \emph{bje }\textless{} *[b]raj (18-17a, 18-16d) `fatigue',
Tib. {\tib{བརྒྱལ་}} \emph{brgyal} \textless{} *bryal (Li's law cf. §) `sink
down, faint'

\item
Chi. 蝸  \emph{kwae} \textless{} *{\sil{kˤ}}roj (19-04c) `snail', WBur. {\burm{ကြွေ
}} \emph{krwe} \textless{} *kruy `shell'

\item
Chi. 橢  \emph{thwaX} \textless{} *l̥{\sil{ˤ}}ojʔ (19-16c) `oval', WBur. {\burm{လွှား
}} \emph{lhwāḥ} \textless{} *hlōḥ `oblong (shield)'

\item
Chi. 墮  \emph{xjwie} \textless{} *l̥oj (19-16e) `destroy', WBur. {\burm{လွှား
}} \emph{lhwāḥ} \textless{} *hlōḥ `hurdle over, unfurl'

\item
Chi. 唾  \emph{thwaH} \textless{} *t{\sil{ʰ}}{\sil{ˤ}}oj-s (19-17m) `spit', Tib. {\tib{ཐོ་ལེ་
}} \emph{tho-le} `spit', WBur. {\burm{ထွေး}} \emph{thweḥ} \textless{} *thuyḥ
'spittle'

\item
Chi. 氐  \emph{tejX} \textless{} *t{\sil{ˤ}}ijʔ (26-14a) `bottom', Tib. {\tib{མཐིལ་
}} \emph{mthil} `bottom, base'

\item
Chi. 虺  \emph{{\sil{xjwɨjX}}} \textless{} *m̥rujʔ (27-19a) `snake' (see p. n. ),
Tib. {\tib{སྦྲུལ་}} \emph{sbrul} \textless{} *smrul (Simon's law, cf. §), OBur. {\burm{မြုယ်}} \emph{mruy}

\item
Chi. 違 \emph{h{\sil{jwɨj}}} \textless{} *{\sil{[ɢ]ʷ}}ə{[}j] (28-05d) `go
against', Tib. {\protect\hypertarget{anchor-60}{}{}} Chi. 違 \emph{h{\sil{jwɨj}}}
\textless{} *[ɢ]{\sil{ʷ}}ə[j] (28-05d) `go against', Tib. {\tib{འགོལ་}}
{\emph{ḫgol}} \textless{} *ḫg{\sil{ʷ}}al `part, deviate'

\item
Chi. 壞 \emph{hweajH} \textless{} *N-{[}k]{\sil{ˤ}}ruj-s (28-06d) `destroy,
ruin', Tib. {\tib{འདྲུལ་}} \emph{ḫdrul} `rot (v.)', Tib. {\tib{བྲུལ་}}~ \emph{brul}
'crumbles'
\end{description}

\paragraph{*-r \textless{} *-r, *-rl}. The correspondence of
Chinese *-r to Tibetan -\emph{l} is difficult to reconstruct with
confidence (see Hill 2014 `cognates'). A mechanical solution is to
reconstruct *-rl which changes to *-r in Chinese and -\emph{l} in
Tibetan (cf. §). Although Auslaut /rl/ is rare in Asia, in many
languages, like rhotic varieties of English, it is plentiful (e.g.
'curl', `furl', `pearl'). Burmese loses both *-r and *-l, so it
naturally loses *-rl also (cf. §), but keeps it as -\emph{y} after the
vowel -\emph{u}- (cf. §). In evaluating the coming comparisons, one must
bear in mind that in Baxter \& Sagart's reconstruction *-[n]
explicitly encompasses *-r as an interpretation.

§a. Examples to be reconstructed *-r :

\begin{description} \item
Chi. 裹 \emph{kwaX} \textless{} *s.{[}k]{\sil{ˤ}}o{[}r]ʔ (19-02d) `wrap
(v.)', Tib. {\tib{སྐོར་}} \emph{skor} `go around'

\item
Chi. 鮮 \emph{sjen} \textless{} *[s]{[}e]r (23-21a) `fresh', Tib. {\tib{གསར་}} \emph{gsar} `new', Bur. {\burm{သ}} \emph{sa} `titivate'

\item
Chi. 竿 \emph{kan} \textless{} *{\sil{kˤ}}ar (24-01k) `pole, rod', Tib. {\tib{མཁར་
}} \emph{mkhar} / {\tib{འཁར་}} \emph{ḫkhar} `staff, stick'

\item
Chi. 獻 \emph{xjonH} \textless{} *ŋ̊ars (24-17e) `offer, present, wise
man', Chi. 義 \emph{ngjeH} \textless{}~*ŋ(r)ajs (18-05r) `duty,
justice', Tib. {\tib{སྔར་}} \emph{sṅar} `intelligent, quick of apprehension'

\item
Chi. 癉 \emph{tanX} \textless{} *t{\sil{ˤ}}arʔ (24-21l) `disease, suffering,
distress', Tib. {\tib{ལྡར་}} \emph{ldar} `be weary, tired, faint'

\item
Chi. 纏 \emph{drjen }\textless{} *[d]ra[n] (24-28c) `bind,
wind', Tib. {\tib{སྟར་}} \emph{star} `tie fast, fasten to'

\item
Chi. 難 \emph{nan} \textless{} *n{\sil{ˤ}}ar (24-35d) `difficult', Tib. {\tib{མནར་
}} \emph{mnar} `suffer, be tormented' (cf. p. n. ), Bur. {\burm{နာ}} \emph{nā}
'hurt'

\item
Chi. 霰 \emph{senH} \textless{} *[s]{\sil{ˤ}}e[n]-s (24-44d) `sleet',
Tib. {\tib{སེར་}} \emph{ser} `hail'

\item
Chi. 緩 \emph{hwanX }\textless{} *{\sil{[ɢ]ʷˤ}}a[n]ʔ (25-14l) `slack;
slow', Tib. {\tib{འགོར་}} \emph{ḫgor} \textless{} *ḫg{\sil{ʷ}}ar `tarry, linger'

\item
Chi. 飛 \emph{p{\sil{jɨj}}} \textless{} *Cə.pə{[}r] (27-09a) `fly (v.)', Chi.
��翂 \emph{pjun} \textless{} *(Cə.)pə{[}r] (33-30ef) `fly (v.), soar',
Chi. 奮 \emph{pjunH} \textless{} *p{[}ə][n]-s (33-33a) `spread
wings and fly', Tib. {\tib{འཕུར་}} \emph{ḫphur} `fly (v.)'

\item
Chi. 銑 \emph{senX} \textless{} *sərʔ (33-25h) `glossy', Tib. {\tib{གསེར་
}} \emph{gser} `gold'
\end{description}

§b. Examples to be reconstructed *-rl :

\begin{description} \item
Chi. 遣 \emph{khjenX} \textless{} *[k]{\sil{ʰ}}e[n]ʔ (23-04b) `send
away', Tib. √skyal (pres. {\tib{སྐྱེལ་}} \emph{skyel}) `send'

\item
Chi. 運 \emph{hjunH} \textless{} *{\sil{[ɢ]ʷ}}ər-s (23-13d) `move', Tib. {\tib{འགུལ་}} \emph{ḫgul} `move'

\item
Chi. 徧 \emph{penH} \textless{} *p{\sil{ˤ}}e[n]-s (23-27b) `(go) all
around', Tib. √pel (pres. {\tib{འཕེལ་}}) `increase, augment'

\item
Chi. 扞捍 \emph{hanH} \textless{} *m-{\sil{kˤ}}a{[}r]-s (24-01q, 24-01i')
'shield (n.), ward off', Chi. 干 \emph{kan} \textless{} *{\sil{kˤ}}a{[}r]
(24-01a) `protect, guard', Tib. {\tib{འགལ་}} \emph{ḫgal} `oppose, contradict',
Bur. {\burm{ကာ}} \emph{kā} `shield (n.)'

\item
Chi. 肝 \emph{kan} \textless{} *s.{\sil{kˤ}}a{[}r] (24-01l) `liver', Tib. {\tib{མཁལ་
}} \emph{mkhal} `kidney, reins', Bur. {\burm{ခါး}} \emph{khāḥ} `loins, waist'

Chi. 扞 \emph{hanH }\textless{} *m-{\sil{kˤ}}a{[}r]-s (24-01q) `fend off',
Tib. {\tib{འགལ་}} \emph{ḫgal} `oppose'

\item
Chi. 鼾 \emph{xan} \textless{} *[{\sil{qʰ}}]{\sil{ˤ}}a{[}r]ʔ (24-01-) `snore',
Tib. {\tib{ཧལ་}} \emph{hal} `pant, snort'

\item
Chi. 鞬 \emph{kjon} \textless{} *ka{[}r] (24-08c) `quiver', Tib. {\tib{རྐྱལ་པ་}} \emph{rkyal-pa} `sack, bag'

\item
Chi. 炭 \emph{thanH} \textless{} *[t{\sil{ʰ}}]{\sil{ˤ}}a[n]s (24-24a) `charcoal,
coal', Tib. {\tib{ཐལ་}} \emph{thal} `dust, ashes'

\item
Chi. 餐 \emph{tshan} \textless{} *ts{\sil{ʰ}}{\sil{ˤ}}ar (24-40c) `eat, food, meal',
Tib. {\tib{མཚལ་མ་}} \emph{tshal-ma} `breakfast'

\item
Chi. 蹯 \emph{bjon} \textless{} *bar (24-54l) `paw', Tib. {\tib{ཕྱག་སྦལ་
}} \emph{phyag}-\emph{sbal} `soft part of an animal's paw'

\item
Chi. 蕃 \emph{bjon} \textless{} *[b]ar (24-54m) `ample, flourish',
Tib. {\tib{དཔལ་}} \emph{dpal} `glory', Tib. {\tib{སྦལ་མིག་}} \emph{sbal-mig} `bud,
sprout'

\item
Chi. 援 \emph{hjwon} \textless{} *{\sil{[ɢ]ʷ}}a[n] (25-14e) `pull up',
Tib. {\tib{འགྲོལ་}} \emph{ḫgrol} \textless{} *ḫg{\sil{ʷ}}ral `become free'

\item
Chi. 西 \emph{sej} \textless{} *s-n{\sil{ˤ}}ər (26-31a) `west' (cf. Sagart 2004:
71-74), Tib. {\tib{མནལ་}} \emph{mnal} `sleep', Bur. {\burm{နား}} \emph{nāḥ} `rest, stop a
while'

\item
Chi. 眉 \emph{mij} \textless{} *mrər (27-14a) `eyebrow', WBur. {\burm{မွေး
}} \emph{mweḥ} \textless{} *muyḥ `body hair'

\item
Chi. 根 \emph{kon} \textless{} *[k]{\sil{ˤ}}ə[n] (33-01b) `root, trunk',
Tib. {\tib{ཁུལ་མ་}} \emph{khul-ma} `bottom or side of sth'

\item
Chi. 銀 \emph{ngin} \textless{} *ŋrə[n] (33-01k) `silver,' Tib. {\tib{དངུལ་}} \emph{dṅul}, OBur. {\burm{ငုယ်}} \emph{ṅuy}

\item
Chi. 齦 \emph{ng{\sil{jɨ}}n} \textless{} *ŋə[n] (33-01-) `gums', Tib.
\protect\hypertarget{anchor-61}{}{}Chi. 齦  \emph{ng{\sil{jɨ}}n} \textless{}
*ŋə[n] (33-01-) `gums', Tib. {\tib{རྙིལ་}} \emph{rñil} \textless{}
*rṅʸil\footnote{The vowel correspondences in this word is irregular; it
  is perhaps for this reason that Jacques rejects this comparison (2013:
  295 note 7).}

\item
Chi. 分 \emph{pjun} \textless{} *pə[n] (33-30a) `divide', Tib. {\tib{འབུལ་
}} \emph{ḫbul}, 
{\tib{འཕུལ་}} \emph{ḫphul} `give'

\item
Chi. 塵 \emph{drin} \textless{} *[d]rə[n] (33-17a) `dust', Tib. {\tib{རྡུལ་}} \emph{rdul}

\item
Chi. 貧 \emph{bin} \textless{} *(Cə.){[}b]rə[n] (33-30v) `poor',
Tib. {\tib{དབུལ་}} \emph{dbul}

\item
Chi. 郡 \emph{gjunH} \textless{} *gur-s (34-12g) `district', Tib. {\tib{ཁུལ་
}} \emph{khul} `district, province'

\item
Chi. 軍 \emph{kjun} \textless{} *[k]{\sil{ʷ}}ər (34-13a) `army', Tib. {\tib{གཡུལ་
}} \emph{g.yul} `army, battle'

\item
Chi. 煇輝 \emph{{\sil{xjwɨj}}} \textless{} {\sil{*qʷʰ}}ər (34-13k, 34-13l) `brilliant',
Tib. {\tib{ཁྲོལ་ཁྲོལ་}} \emph{khrol-khrol} \textless{} *kh{\sil{ʷ}}ral `bright, shining,
sparkling, glistening'
\end{description}

§c. It is not always possible to confidently assign a cognate set to
Trans-Himalayan *-r or *rl, as the following example shows.

\begin{description} \item
Chi. 䫀 \emph{konX} \textless{} *[k]{\sil{ˤ}}ə[n]ʔ (33-01-) `neck', Tib. {\tib{མགུལ་}} \emph{mgul} `neck', Tib. 
%{\protect\hypertarget{anchor-62}{}{}Chi.
%䫀 \emph{konX} \textless{} *[k]{\sil{ˤ}}ə[n]ʔ (33-01-) `neck', Tib. 
{\tib{མགུལ་}} \emph{mgul} `neck', Tib. {\tib{མགུར་}} \emph{mgur} `neck'

\item
Chi. 粲 \emph{tshanH} \textless{} *[ts{\sil{ʰ}}]{\sil{ˤ}}ars (24-40b) `bright and
white', Tib. {\tib{མཚར་}} \emph{mtshar} `fair, beautiful, bright', Tib. √stsal
(pres. 
%{\protect\[hypertarget{anchor-63}{}{}
Chi. 粲 \emph{tshanH}
\textless{} *[ts{\sil{ʰ}}]{\sil{ˤ}}ars (24-40b) `bright and white', Tib. {\tib{མཚར་}} \emph{mtshar} `fair, beautiful, bright', Tib. √stsal (pres. {\tib{གསལ་
}} \emph{gsal}, cf. Hill 2012 `six vowels' : 25) `clear, bright'
\end{description}

\paragraph{The correspondence of Chinese -o- to Tibetan -u-:} Because
there are examples in which both Tibetan and Chinese have -u- and
examples in which both languages have -o-, those examples of Tibetan
-\emph{u}- cor­re­spond­ing to Chinese *-o- cannot be reconstructed as
either *u or *o. Consequently it is necessary to reconstruct a vowel
unattested in both languages, which changes into -\emph{u}- in Tibetan
and *-o- in Chinese. I propose to reconstruct this correspondence as
*-əw-, largely because this fills a gap in Old Chinese (Hill 2012
'evolution': 75-77, 2012 `six vowels': 32-35, cf. §). This is a
tentative suggestion, which faces two potential objections. First, it is
worrisome that examples of *-əw- outnumber those of *-o-, because
 \emph{a priori} *-əw- should be less common than *-o- in the
proto-language. Second, if *-aw and *-ew merge to -\emph{o} in Tibetan
(cf. §), one might expect *-əw- to also yield -\emph{o} in Tibetan.
However, the fact that this re­con­struc­tion is called for only in open
syllables or syllables with velar codas (with 洽 \emph{heap} \textless{}
*{\sil{ɢˤ}}rop as the one exception), by pa­ral­le­ling the distribution of *-w
in Old Chinese argues in favour of this reconstruction. The Burmese
cognates, omitted here, are treated in section §.

§a. Examples to be reconstructed *-u-:

\begin{description} \item
Chi. 胞 \emph{paew} \textless{} *p{\sil{ˤ}}ru (13-72b) `womb', Tib. {\tib{ཕྲུ་མ་}} \emph{phru-ma} `afterbirth'
 \item
Chi. 舅 \emph{gjuwX} \textless{} *[g](r)uʔ (04-16b) `maternal
uncle', Tib. {\tib{ཁུ་}} \emph{khu} `paternal uncle'
\item
Chi. 九 \emph{kjuwX} \textless{} *[k]uʔ (04-12a) `nine', Tib. {\tib{དགུ་
}} \emph{dgu}
\item
Chi. 鳩 \emph{kjuw} \textless{} *[k](r)u (04-12n) (a kind of bird),
Tib. {\tib{འང་གུ་}} \emph{ḫaṅ-gu} `pigeon'
\item
Chi. 嗥 \emph{haw} \textless{} *g{\sil{ˤ}}u (13-01d) `roar, wail', Tib. {\tib{ངུ་}} \emph{ṅu} `weep'
\item
Chi. 肘 {\emph{trjuwX}} \textless{} *t.kruʔ (13-23a) `elbow', Tib. {\tib{གྲུ་མོ་}} \emph{gru-mo}
\item
Chi. 流 \emph{ljuw} \textless{} *[r]u (13-46a) `flow', Tib. {\tib{རྒྱུ་
}} \emph{rgyu} \textless{} *ryu (Li's law, cf. §)
\item
Chi. 篤  \emph{towk} \textless{} *t{\sil{ˤ}}uk (14-08g) `firm, solid', Tib. {\tib{འཐུག་}} \emph{ḫthug}, {\tib{མཐུག་}} \emph{mthug} `thick, dense'
\item
Chi. 晝  \emph{trjuwH} \textless{} *truks (14-09a) `time of daylight',
Tib. {\tib{གདུགས་}} \emph{gdugs} `mid-day, noon'
\item
Chi. 覺  \emph{kaewk} \textless{} *{\sil{kˤ}}ruk (14-03f) `awake', Tib. {\tib{དཀྲུག་
}} \emph{dkrug} `stir, agitate, disturb'
\end{description}

A complete list of examples of this correspondence is given at §.

§b. Examples to be reconstructed *-o-:

\begin{description} \item
Chi. 絶  \emph{dzjwet} \textless{} *[dz]ot (22-16a) `cut off, break
off', Tib. {\tib{ཆོད་}} \emph{chod} `be sharp', WBur. {\burm{ဆွတ်}}~ \emph{chwat}
\textless{} *chot `pluck'
\item
Chi. 脫  \emph{thwat} \textless{} *mə-l̥{\sil{ˤ}}ot (22-13m) `peel off', Tib.
{\tib{གློད་}} \emph{glod} `loose, relaxed', OBur. {\burm{လောတ်~}} \emph{lo{\sil{₁}}t} `be free'
\item
Chi. 倌  \emph{kwaenH} \textless{} *kr{\sil{ˤ}}ons (25-01l) `servant, groom',
Tib. {\tib{ཁོལ་}} \emph{khol} `servant', OBur. {\burm{ကျောန်}} \emph{kyo{\sil{₁}}n} `slave'
\item
Chi. 垂  \emph{dzywe} \textless{} *[d]oj (19-17a) `hang down', Tib. {\tib{འཇོལ་}} \emph{ḫǰol} `hang down', WBur. {\burm{ဆွဲ}} \emph{chwai }\textless{} *choi
\item
Chi. 巷 \emph{haewngH} \textless{} *C.{[}g]{\sil{ˤ}}roŋ-s `lane, street', Tib. {\tib{གྲོང་}} \emph{groṅ} `village'
\item
Chi. 卵  \emph{lwanX} \textless{} *k.r{\sil{ˤ}}orʔ (25-32a) `egg', Tib. {\tib{སྲོ་མ་}} \emph{sro-ma} `louse egg'
\item
Chi. 殼  \emph{khaewk} \textless{} *[{\sil{kʰ}}]{\sil{ˤ}}rok (11-03a) `hollow shell,
hollow', Tib. {\tib{སྐོག་}} \emph{skog} `shell, peel',
\item
Chi. 蔥 \emph{tshuwng} \textless{} *[ts]{\sil{ʰ}}{\sil{ˤ}}oŋ (12-19g) `onion', Tib. {\tib{བཙོང་}} \emph{btsoṅ}
\item
Chi. 綴 \emph{trjwet} \textless{} *trot (22-10b) `bind', Chi. 贅
 \emph{tsywejH} \textless{} *tots (22-11a) `unite, together', Tib. √rtod
(pres. {\tib{རྟོད་}}) `tether, fasten, secure'
\item
Chi. 悅  \emph{ywet} \textless{} *lot (22-13o) `pleased', Tib. {\tib{བྲོད་}} \emph{brod} `joy, joyful'
\item
Chi. 掘  \emph{gjwot} \textless{} *[g]ot (31-16s) `dig out (earth)',
Tib. {\tib{རྐོ་}} \emph{rko} `dig'
\item
Chi. 涫  \emph{kwanH} \textless{} *{\sil{kˤ}}ons (25-01f) `bubble', Tib. {\tib{འཁོལ་}} \emph{ḫkhol} `boil'
\item
Chi. 唾  \emph{thwaH} \textless{} *t{\sil{ʰ}}{\sil{ˤ}}ojs (19-17m) `spit', Tib. {\tib{ཐོ་ལེ་}} \emph{tho-le} `spit'
\item
Chi. 鑽  \emph{tswan} \textless{} *[ts]{\sil{ˤ}}or (24-39h) `perforate,
penetrate', Chi. 鐫 \emph{tsjwen} \textless{} *tson (25-39c) `chisel,
sharp point', Tib. {\tib{མཚོན་}} \emph{mtshon} `weapon'
\item
Chi. 裹  \emph{kwaX} \textless{} *s.{[}k]{\sil{ˤ}}o{[}r]ʔ (19-02d) `wrap
(v.)', Tib. {\tib{སྐོར་}} \emph{skor} `go around'
\end{description}

§c. Examples to be reconstructed *-əw-:

\begin{description} \item
Chi. 寇  \emph{khuwH} \textless{} *[k]{\sil{ʰ}}{\sil{ˤ}}(r)os (10-04a) `steal', Tib. {\tib{རྐུ་}} \emph{rku}
\item
Chi. 曲  \emph{khjowk} \textless{} *{\sil{kʰ}}(r)ok (11-04a) `bent, crooked',
Tib. {\tib{འགུགས་}} \emph{ḫgugs} `bend'
\item
Chi. 觸  \emph{tsyhowk} \textless{} *t{\sil{ʰ}}ok (11-12g) `knock against', Tib. {\tib{གཏུག་}} \emph{gtug} `meet, touch'
\item
Chi. 乳  \emph{nyuX} \textless{} *noʔ (10-32a) `milk, nipple', Tib. {\tib{ནུ་}} \emph{nu} `suck'
\item
Chi. 孺  \emph{nyuH}\textless{} *nos (10-31d) `child, mild', Tib. {\tib{ནུ་བོ་}} \emph{nu-bo} `younger brother'
\item
Chi. 霧  \emph{mjuH} \textless{} *kə.m(r){[}o]ks (13-76t) `fog, mist',
Tib. {\tib{རྨུགས་}} \emph{rmugs} `dense fog'
\item
Chi. 撞  \emph{draewng} \textless{} *[N-t]{\sil{ˤ}}roŋ (12-08f') `strike',
Tib. {\tib{རྡུང་}} \emph{rduṅ} `strike, beat'
\item
Chi. 蜂�� 
\emph{phjowng} \textless{} *p{\sil{ʰ}}(r)oŋ (12-25st) `bee', Tib. {\tib{བུང་བ་}} \emph{buṅ-ba}
\item
Chi. 空  \emph{khuwng} \textless{} *{\sil{kʰ}}{\sil{ˤ}}oŋ (12-01h) `hollow, empty, hole',
Chi. 孔  \emph{khuwngX} \textless{} *{\sil{kʰ}}oŋʔ (12-02a) `empty', Tib. {\tib{ཁུང་}} \emph{khuṅ} `hole, pit, hollow, cavity'
\item
Chi. 洽  \emph{heap} \textless{} *{\sil{[ɢ]ˤ}}r{[}o]p (37-01m) `accord
with', Tib. {\tib{འགྲུབ་}} \emph{ḫgrub} `accomplish, achieve'
\end{description}

A complete list of examples of this correspondence is given at §.

\section{Reprise: Trans-Himalayan to Old
Chinese}\label{reprise-trans-himalayan-to-old-chinese}

\paragraph{}Since the changes from Trans-Himalayan to Old Chinese do not interact
with one another it is not possible to present them in chronological
order. Nonetheless, for the sake of consistency the reprise presents
them in the opposite order to their more detailed presentation above.

\begin{description} \item
*əw \textgreater{} *-o- (cf. §)

\item
*-r, *-rl \textgreater{} *-r (cf. §)

\item
*-j, *-l \textgreater{} *-j (cf. §)

\item
*-ʔ, *-q \textgreater{} *-ʔ (cf. §)

\item
*-k, *-kə \textgreater{} *-k (cf. §)

\item
Pulleyblank's conjecture: *RC- \textgreater{} *Cr- (cf. §)

\item
Coblin's conjecture: *rT- \textgreater{} *Tr- (cf. §)

\item
Inability to recognize *-r- in certain circumstances (cf. §)

\item
Inability to distinguish *-u- and in -ə- before acutes (cf. §)
\end{description}

\begin{description} \item
Labial neutralization (cf. §)

\item
*-iŋ, *-in \textgreater{} *-in and *-ik, *-it \textgreater{} -it (cf. §)
\end{description}

\section{Diachronic mysteries}\label{diachronic-mysteries-2}

\paragraph{}The most valuable contribution of a survey of Chinese sound laws is
to draw new focus to the exceptions to these sound laws. After having
surveyed what is currently known of Chinese historical phonology, those
areas in need of better study merit focus. Exceptions to the respective
sound laws presented have been provided above, but it is convenient to
assemble them together here. The word 汝  \emph{nyoX} \textless{} *naʔ
(01-56j) `you' does not regularly correspond to Burmese 
{\burm{နင်}} \emph{naṅ}
'you', but it seems likely the two are cognate (cf. §c). The word 膚
 \emph{pju} \textless{} *pra (01-51g) `skin' is irregularly missing a
final glottal stop (cf. §c and p. n. ). The word 沙  \emph{srae}
\textless{} *s{\sil{ˤ}}raj (18-15a) `sand' is an exception to Pulleyblank's
conjecture (cf. §a); the medial *-r- in this word is unexpected from a
comparative context.
